\documentclass[11pt,a4paper]{article}
\usepackage[T1]{fontenc}
\usepackage[ngerman]{babel}
\usepackage[hmargin=29mm,vmargin=25mm]{geometry}
\usepackage{amsmath,amssymb}
\usepackage{tikz}
\usepackage{fvextra}
\usetikzlibrary{arrows.meta}
\definecolor{qs ink}{HTML}{242B33}
\definecolor{qs axis}{HTML}{565F68}
\definecolor{qs blue}{HTML}{174D75}
\definecolor{qs rust}{HTML}{99471C}
\usepackage[hidelinks,pdftitle={Quantum Sphaera},pdfauthor={Marc Brendecke},pdfsubject={Überarbeitungsentwurf zu Preprint 0.1.0}]{hyperref}

\setlength{\parindent}{1em}
\raggedbottom

\newcommand{\Tr}{\operatorname{tr}}
\newcommand{\id}{\mathrm{I}}
\begin{document}
\thispagestyle{plain}
\noindent{\LARGE\bfseries Quantum Sphaera\par}
\vspace{0.25em}
\noindent{\large Auf dem Weg zu einem gemeinsamen Rahmen\\
für Materie, Eichstruktur und Raumzeitgeometrie\par}
\vspace{0.7em}
\noindent Marc Brendecke\hfill{\small Überarbeitung zu 0.1.0 \enspace$\cdot$\enspace 22. November 2026}\par
\noindent{\small Veröffentlichte Fassung \href{https://doi.org/10.5281/zenodo.24405553}{doi:10.5281/zenodo.24405553}\par}
\vspace{0.55em}
\section*{Vorwort}
Was soll ein Heisenberg-Kompensator eigentlich kompensieren? Und wenn man dann mal endlich erfährt, was er kompensieren soll, fragt man sich ziemlich erstaunt: \glqq Wie verdammt noch mal soll er das machen?\grqq{}

Also, was tut ein Star-Trek-Nerd mit zu viel Zeit und einem LLM? Richtig, er sagt ihm: \glqq Hey, erklär mir doch mal, wie das Ding bei \emph{Star Trek} funktionieren soll.\grqq{}

Und Michael Okuda aus dem \emph{Star-Trek}-Team hatte darauf tatsächlich schon eine ziemlich genaue Antwort.

\begin{quote}
\glqq Danke der Nachfrage, er funktioniert sehr gut!\grqq{}\footnote{Freie Übersetzung der Antwort von Michael Okuda in \emph{Time}, 28. November 1994. \href{https://content.time.com/time/subscriber/article/0\%2C33009\%2C981892\%2C00.html}{Reconfigure the Modulators!}}
\end{quote}

Gut, an der Stelle kommt man wohl nicht weiter. Also wird das LLM jetzt etwas gequält. Und der geneigte Leser vielleicht auch. ;)

\noindent Für die mitunter eigenwillige Namensgebung bitte ich um Nachsicht und gehe \mbox{\emph{in medias res}}.

Wieso sollte man Heisenberg überhaupt kompensieren wollen, wenn man beamen möchte?

Und, was passiert mit unserer Physik, wenn wir einfach mal annehmen, dass es so was wie einen Heisenberg-Kompensator wirklich geben könnte?
\medskip

\section{Vom Photon zum Kompensator}
\label{sec:question}

Angefangen hat das Ganze für mich beim Photon.

Was muss eine Konstruktion beschreiben können, damit sie als Modell für Licht infrage kommt? Unser erster Prüfstein ist eine ebene elektromagnetische Welle im Vakuum. In Einheiten mit \(c=1\) sind ihre Feldvektoren gleich lang und orthogonal. An diesen beiden Eigenschaften prüfen wir die geplante Operation. Die Frage nach einem quantisierten Lichtzustand schließt sich an diesen klassischen Vergleich an.

Im Vorwort stand die Frage, was ein Heisenberg-Kompensator für die Physik bedeuten könnte. Der erste Schritt ist kleiner. Wir suchen eine Operation, die einen bestimmten Anteil eines Zustands aufhebt und den Rest erhält. Wie kann sie die beiden Anteile unterscheiden?

\subsection{Was der Kompensator tun soll}
\label{sec:carrier}

Stellen wir uns einen Zustand vor, den wir einmal spiegeln. Ein Teil bleibt dabei gleich, ein anderer wechselt das Vorzeichen. Legen wir den ursprünglichen und den gespiegelten Zustand übereinander und bilden den Mittelwert, fällt der zweite Teil heraus. Damit ist die Aufgabe des Kompensators beschrieben. Jetzt brauchen wir eine Regel für die Spiegelung.

Sei \(J:V\to V\) eine lineare Involution auf einem reellen Vektorraum \(V\). Um zu sehen, was der Mittelwert auswählt, schreiben wir \(v=v_++v_-\) mit \(Jv_+=v_+\) und \(Jv_-=-v_-\). Dann ist \(Jv=v_+-v_-\). Addieren und Subtrahieren liefert den gesuchten Anteil und sein Gegenstück:
\begin{equation}
 Cv=\frac{v+Jv}{2}=v_+,\qquad
 P_-v=\frac{v-Jv}{2}=v_-.
 \label{eq:projections}
\end{equation}
Umgekehrt erfüllen diese beiden Ausdrücke für jedes \(v\) die angegebenen Vorzeichenregeln. Damit bestimmt \(J\) die beiden Anteile eindeutig. Aus \(J^2=I\) folgen \(C^2=C\), \(P_-^2=P_-\), \(CP_-=0\) und \(C+P_-=I\). Der erste Operator ist der Heisenberg-Kompensator, kurz HC, in der hier betrachteten Realisierung \cite{physics}. Wenn der Raum ein Skalarprodukt trägt und \(J\) orthogonal ist, sind auch die beiden ausgewählten Räume orthogonal. Die Zerlegung selbst folgt schon aus der Involution.

Der Name greift das Motiv der quantenmechanischen Unschärfe auf. Die Heisenberg-Unschärferelation betrifft die Streuungen von Observablen wie Ort und Impuls \cite{robertson}. Was hier kompensiert wird, ist vorerst genau der Anteil, den \(J\) umkehrt. Eine physikalische Wirkung lässt sich erst beurteilen, sobald wir wissen, was \(v\) beschreibt und welche Operation \(J\) auf diesen Daten ausführt.

\subsection{Was der Kompensator am Lichtfeld auswählt}
\label{sec:field-compensation}

Kehren wir zur Lichtwelle zurück. In einer Normierung mit \(c=1\) schreiben wir die Felder als \(\mathbf z=\mathbf E+i\mathbf B\). Ein einfacher Kandidat für \(J\) ist der Austausch von elektrischem und magnetischem Feld. In dieser Schreibweise lautet er \(J_{\mathrm B}\mathbf z=i\overline{\mathbf z}\). Das ist eine reell-lineare Involution, mit \(J_{\mathrm B}(\mathbf E+i\mathbf B)=\mathbf B+i\mathbf E\). Setzen wir sie in \eqref{eq:projections} ein, erhalten wir
\begin{equation}
 C\mathbf z=(1+i)\frac{\mathbf E+\mathbf B}{2},
 \qquad
 P_-\mathbf z=(1-i)\frac{\mathbf E-\mathbf B}{2}.
 \label{eq:hc-field}
\end{equation}
Der HC gleicht die Differenz aus. Im ausgewählten Anteil sind beide Feldvektoren gleich und haben den Wert \((\mathbf E+\mathbf B)/2\) \cite{grundidee,feldsektoren}.

\paragraph{Eine ebene Welle als Eingabe.}
Nehmen wir an einem Punkt der Welle \(\mathbf E=a\mathbf e_x\) und \(\mathbf B=a\mathbf e_y\), mit \(a>0\) und orthonormalen Raumrichtungen \(\mathbf e_x,\mathbf e_y\). Die Welle hat dort den Feldvektor \(\mathbf z=a(\mathbf e_x+i\mathbf e_y)\). Gleichung \eqref{eq:hc-field} liefert
\begin{equation}
 C\mathbf z=\frac a2(1+i)(\mathbf e_x+\mathbf e_y),\qquad
 P_-\mathbf z=\frac a2(1-i)(\mathbf e_x-\mathbf e_y).
 \label{eq:field-example}
\end{equation}
Im ersten Anteil sind beide Feldvektoren gleich. Im zweiten sind sie entgegengesetzt. Ihre Summe rekonstruiert die Eingabe. Im euklidischen Feldskalarprodukt ist das ursprüngliche Quadrat der Norm \(2a^2\), und jeder der beiden Anteile hat die quadrierte Norm \(a^2\). Das ist die Pythagoras-Identität der reellen orthogonalen Zerlegung. Für eine Aussage über Energietransport wäre zusätzlich die Feldentwicklung zu prüfen \cite{feldsektoren,grundidee}.

\begin{figure}[!htbp]
\centering
\begin{tikzpicture}[>=Latex,x=2.35cm,y=2.35cm, font=\small]
  \begin{scope}
    \draw[qs axis,->] (-0.12,0)--(1.25,0);
    \draw[qs axis,->] (0,-0.12)--(0,1.22);
    \draw[qs blue,very thick,->] (0,0)--(1,0) node[below=4pt] {\(\mathbf E\)};
    \draw[qs rust,very thick,->] (0,0)--(0,1) node[left=4pt] {\(\mathbf B\)};
    \draw[qs axis] (0.11,0)--(.11,.11)--(0,.11);
    \node[font=\small\bfseries,anchor=south west] at (-.17,1.34) {Eingang};
    \node[anchor=north west,align=left] at (-.17,-.25)
      {\(|\mathbf E|=|\mathbf B|\), \(\mathbf E\cdot\mathbf B=0\)\\\(s=0\) bei nichtverschwindendem Feld};
  \end{scope}
  \draw[->,thick,qs ink] (1.59,.43)--(2.35,.43)
     node[midway,above=4pt] {\(C\)};
  \begin{scope}[xshift=6.45cm]
    \draw[qs axis,->] (-.12,0)--(1.25,0);
    \draw[qs axis,->] (0,-.12)--(0,1.22);
    \draw[qs axis,densely dotted] (.5,0)--(.5,.5)--(0,.5);
    \draw[qs blue,very thick,->] (0,0)--(.5,.5);
    \node[anchor=south west,align=left] at (.51,.5) {\(\mathbf E'=\mathbf B'\)};
    \node[font=\small\bfseries,anchor=south west] at (-.17,1.34) {Nach der Projektion};
    \node[anchor=north west,align=left] at (-.17,-.25)
      {\(\mathbf E'=\mathbf B'=(\mathbf E+\mathbf B)/2\)\\\(s'=2i|\mathbf E'|^2\)};
  \end{scope}
\end{tikzpicture}
\caption{Ein lichtartiges Eingangsfeld und sein Bild unter der HC-Projektion. Die Projektion mittelt die beiden Feldvektoren. In der Zeichnung haben die orthogonalen Eingangsvektoren gleiche Länge. Ihre beiden Ausgänge liegen auf demselben Pfeil.}
\label{fig:hc-feld}
\end{figure}


Und ist der ausgewählte Anteil jetzt Licht? Aus \eqref{eq:field-example} folgen für die beiden projizierten Feldvektoren \(|\mathbf E'|^2=|\mathbf B'|^2=a^2/2\) und \(\mathbf E'\cdot\mathbf B'=a^2/2\). Die Projektion erhält in diesem Beispiel also die Gleichheit der Beträge. Die Orthogonalität geht verloren. Ein Kennzeichen für lichtartige Felder muss beide Bedingungen zugleich prüfen.

\subsection{Wie die Sektoren ins Spiel kommen}
\label{sec:frames}

Die beiden Eigenschaften der Welle lassen sich in einer Zahl zusammenfassen. Mit dem komplex-bilinearen Skalarprodukt ist
\begin{equation}
 s(\mathbf z)=\mathbf z\cdot\mathbf z
 =|\mathbf E|^2-|\mathbf B|^2+2i\,\mathbf E\cdot\mathbf B .
 \label{eq:field-invariant}
\end{equation}
Verschwinden Real- und Imaginärteil, sind die beiden Felder gleich stark und stehen senkrecht aufeinander. Mit \(\mathbf z\ne0\) haben wir dann ein lichtartiges Feld. Diese Kennzeichnung bleibt bei Lorentztransformationen des Feldes erhalten. Für das Bild aus \eqref{eq:hc-field} mit \(\mathbf E'=\mathbf B'\ne0\) ergibt sich \(s=2i|\mathbf E'|^2\). So lässt sich der Befund direkt an derselben Feldgröße ablesen \cite{grundidee,feldkern}.

Hier kommt die Sektorlesart ins Spiel. Auf dem Schnitt mit reellem \(s\) gilt \(\mathbf E\cdot\mathbf B=0\). Positives \(s\) kennzeichnet dort elektrische Dominanz, negatives \(s\) magnetische Dominanz, und \(s=0\) enthält die lichtartigen Felder sowie das Nullfeld. Außerhalb dieses Schnitts kann \(s\) komplex sein. Die Sektoren bezeichnen hier Klassen von Feldkonfigurationen. Die beiden Anteile der HC-Zerlegung aus \eqref{eq:projections} beantworten eine andere Frage, nämlich ob \(J\) einen Anteil erhält oder umkehrt.

Eine Dualitätsvierteldrehung \(\mathbf z\mapsto i\mathbf z\) wechselt das Vorzeichen von \(s\). Sie verbindet damit elektrische und magnetische Dominanz bei reellem \(s\). Der Betrag \(|s|\) bleibt dabei erhalten. Ein Feld mit \(s\ne0\) bleibt daher auch nach dieser Drehung außerhalb der lichtartigen Klasse. Im Companion heißt diese Abbildung Tachyon-Abbildung (Runde v4.10.0-r15 \cite{companion}). Sie ist nicht mit der Zeitumkehr der Feldamplituden zu verwechseln. Diese wirkt als komplexe Konjugation \(\mathbf z\mapsto\overline{\mathbf z}\) (Runde v4.10.0-r31 \cite{layers}). Die Austauschoperation \(J_{\mathrm B}\) aus \eqref{eq:hc-field} enthält zusätzlich die komplexe Konjugation \cite{grundidee,feldsektoren}.

\paragraph{Die Feldklassen in einer gemeinsamen Darstellung.}
Schreiben wir das Felddublett als reelle Matrix \(F=[\,\mathbf E\mid\mathbf B\,]\). In diesen Koordinaten wirkt \(J_{\mathrm B}\) durch Vertauschen der Spalten. Die Gram-Matrix lautet
\begin{align}
 G_F=F^{\mathsf T}F
 &=\begin{pmatrix}
 |\mathbf E|^2 & \mathbf E\cdot\mathbf B\\
 \mathbf E\cdot\mathbf B & |\mathbf B|^2
 \end{pmatrix}, \qquad u_F=\tfrac12\operatorname{tr}G_F,\notag\\
 G_F&=u_F I_2+\tfrac12\bigl(\operatorname{Re}s\,\sigma_z+
                           \operatorname{Im}s\,\sigma_x\bigr).
 \label{eq:field-gram}
\end{align}
Die Pauli-Matrizen \(\sigma_x,\sigma_z\) beziehen sich hier auf die beiden Feldspalten. Für \(u_F>0\) kann man die dimensionslose Auslese \(\mathbf b=(\operatorname{Re}s,\operatorname{Im}s)/(2u_F)\) verwenden. Aus \(\det G_F=|\mathbf E\times\mathbf B|^2\) folgt
\begin{equation}
 |\mathbf b|^2=1-\frac{\det G_F}{u_F^2}\leq1 .
 \label{eq:field-disk-gram}
\end{equation}
Die nichtverschwindende Eingangswelle aus \eqref{eq:field-example} liegt im Zentrum dieser Scheibe, mit \(G_F=a^2 I_2\). Ihre beiden projizierten Anteile liegen bei \(\mathbf b=(0,1)\) und \((0,-1)\). Ihre Gram-Matrizen haben Rang eins, weil die Spalten jeweils parallel sind. So zeigt dasselbe Beispiel, wo die lichtartige Feldklasse liegt und wohin sie der HC abbildet \cite{feldsektoren}.

Die Unterscheidung zwischen den Operationen lässt sich ebenfalls an \(s\) prüfen. Bei einer Dualitätsdrehung um \(\alpha\) gilt \(s(e^{i\alpha}\mathbf z)=e^{2i\alpha}s(\mathbf z)\). Beim Austausch gilt \(s(J_{\mathrm B}\mathbf z)=-\overline{s(\mathbf z)}\). Die Lorentzwirkung auf den komplexen Feldkoordinaten erfolgt durch \(\Lambda\in SO(3,\mathbb C)\) und erhält \(s\) wegen \(\Lambda^{\mathsf T}\Lambda=I_3\). Diese Erhaltung betrifft \(s\) selbst. Die normierte Länge \(|\mathbf b|=|s|/(2u_F)\) hängt zusätzlich vom im jeweiligen Bezug ausgewerteten \(u_F\) ab \cite{feldkern,feldsektoren}.

\paragraph{Was ein Wechsel des Bezugs mit dem HC macht.}
Soll dieselbe Austauschregel \(J_{\mathrm B}\) in zwei Lorentzbezügen gelten, muss sie mit \(\Lambda\) kommutieren. Die beiden Reihenfolgen ergeben
\[
 \Lambda J_{\mathrm B}\mathbf z=i\Lambda\overline{\mathbf z},\qquad
 J_{\mathrm B}\Lambda\mathbf z=i\overline\Lambda\,\overline{\mathbf z}.
\]
Die Ausdrücke stimmen für alle \(\mathbf z\) genau dann überein, wenn \(\Lambda=\overline\Lambda\). Gewöhnliche Raumdrehungen erfüllen das. Ein Boost ändert im Allgemeinen die Form der Austauschregel. Für den transportierten Projektor verwendet man dann \(\Lambda C\Lambda^{-1}\). Die Frage nach einer physikalischen Auslese umfasst daher auch die Wahl und den Transport der Operation \cite{feldsektoren,grundidee}.

Der Companion verwendet für seine physikalischen Modellbezüge die Namen Materieframe \(M\), Photon-Eigenframe \(P\) und Tachyonframe \(T\). Dazu kommt der Bezugsframe \(R\). Für ein nicht lichtartiges Feld ist das der Lorentzbezug, in dem \(\mathbf E\) und \(\mathbf B\) parallel sind. Dort ist \(\det G_F=0\), das Feld liegt also auf dem Rand der Scheibe aus \eqref{eq:field-disk-gram}. Die lichtartigen Felder liegen dagegen in jedem Bezug im Zentrum (Runde v4.10.0-r21 \cite{companion}). Seit Runde v4.10.0-r69 tragen die Frames zusätzlich beschreibende Namen: \(R\) heißt Ruheframe, \(T\) Dual-Frame und \(P\) Lichtpunkt, \(M\) bleibt der Materieframe. Die Buchstaben ändern sich nicht, und dieser Text behält die älteren Namen bei. Ein \emph{Frame} gibt dort an, in welchem Modellbezug ein Zustand dargestellt und eine Größe ausgewertet wird. Eine \emph{Auslese} ist die Abbildung, die aus den Zustandsdaten eine solche Größe bestimmt. So ist \eqref{eq:field-invariant} eine Auslese, mit der wir gerade eine Feldklasse bestimmen. Für den Vergleich zwischen Frames muss für jede relevante Größe angegeben werden, wie sie übertragen wird. Aus dem Feldvergleich folgt die Unterscheidung elektrischer, magnetischer und lichtartiger Konfigurationen. Ihre Identifikation mit den Teilchen und Bezügen des Companion bleibt eine eigene physikalische Aufgabe \cite{companion,grundidee}. Wo die Frames in die offenen Probleme eingehen, fasst Abschnitt~\ref{sec:open-points} zusammen.

Damit können wir die nächste Frage stellen. Welcher gemeinsame Träger kann die Feldzustände, ihre Lorentzwirkung und die algebraischen Operationen aufnehmen und lässt sich zugleich für Materiezustände und ihre Ladungen verwenden?

\subsection{QAM und die Nachvollziehbarkeit der Aussagen}
\label{sec:qam}

QAM (im Companion „Quantum Analysis Mengenlehre“, die Sprach- und Ausleseschicht über der Mengenlehre \cite{companion}) ist die Sprache, mit der der Companion mathematische Objekte und Aussagen registriert. Zu einer Aussage gehören der verwendete Kontext, ihre Voraussetzungen und ihr Status. Sie kann auf eine Herleitung oder auf das Protokoll einer Rechnung verweisen. Ein \emph{Guard} bezeichnet eine noch offene Bedingung, von der ein weitergehender Schluss abhängt. So lässt sich beim Wiederverwenden einer Aussage nachvollziehen, welche Bedingungen für ihre Anwendung gelten \cite{qam}.

Am Kompensator lässt sich das direkt sehen. Die Aussage \(C^2=C\) stützt sich auf die Voraussetzung \(J^2=I\). Für das Ergebnis am elektrischen und magnetischen Feld kommt die Wahl der konkreten Austauschoperation \(J_{\mathrm B}\) hinzu. Ob ein physikalischer Prozess diese Projektion tatsächlich realisiert, bleibt ein Guard. Dieselbe Unterscheidung zwischen Rechnung, Modellannahme und physikalischer Zuordnung wird bei den Leptonenmassen in Abschnitt \ref{sec:mass} wichtig.

Die Formeln stehen im Artikel in ihrer gewöhnlichen mathematischen Schreibweise. In der QAM-Dokumentation können die Werkzeuge die deklarierten Abhängigkeiten und Kontexte prüfen. Die Formel bleibt dabei mathematischer Text. Ihre Herleitung und die physikalische Zuordnung müssen jeweils durch das angegebene Argument oder den Beleg getragen werden.

\section{Was der Träger für Materie und Eichstruktur leisten müsste}
\label{sec:shared-carrier}

\subsection{Die Biquaternionen als Kandidat}

Für die Feldrechnung in Abschnitt \ref{sec:frames} genügte ein komplexer Dreiervektor. Für einen gemeinsamen Träger suchen wir darüber hinaus eine Algebra, in der sich Zustände und Operationen zusammensetzen lassen. Der erste Test für einen solchen Kandidaten ist konkret. Er soll die Feldkoordinaten aufnehmen und ihre Lorentzwirkung samt der Invariante \(s\) wiedergeben. Für die physikalische Erweiterung muss er außerdem Raum für die inneren Symmetrien und die Zustandsauswahl lassen.

Der Companion bettet den Feldvektor in den Vektorteil eines Biquaternions ein. Mit den Quaternion-Einheiten \(e_k\) lautet der Ansatz
\begin{equation}
 q=it\,1+z_1e_1+z_2e_2+z_3e_3\in\mathbb C\otimes_{\mathbb R}\mathbb H,\qquad t\in\mathbb R .
 \label{eq:biquaternion}
\end{equation}
Für die quadratische Norm mit rein quaternionischer Konjugation gilt \(N(q)=-t^2+s(\mathbf z)\). Die Konjugationswirkung durch Biquaternionen der Norm eins trägt auf dem Vektorteil die eigentliche orthochrone Lorentzwirkung und erhält \(s\). Der Skalarparameter \(t\) bleibt unter dieser Wirkung fest. Seine mögliche Deutung ist eine zusätzliche Frage \cite{feldkern,grundidee}.

\paragraph{Welche Invariante der Träger tatsächlich erhält.}
Sei \(q=it\,1+\mathfrak z\) mit \(\mathfrak z=\sum_k z_k e_k\). Aus \(e_j e_k+e_ke_j=-2\delta_{jk}\,1\) folgt \(\mathfrak z^2=-s(\mathbf z)\,1\). Die quadratische Norm \(N(q)\) und ihre Multiplikativität lassen sich damit ausrechnen.
\begin{align}
 q\bar q&=(it\,1+\mathfrak z)(it\,1-\mathfrak z)
       =[-t^2+s(\mathbf z)]\,1,\\
 N(q_1q_2)&=N(q_1)N(q_2).
 \label{eq:biquaternion-check}
\end{align}
Die Produktidentität gilt für \(q_1,q_2\in\mathbb C\otimes_{\mathbb R}\mathbb H\) mit \(N(q)=q\bar q\) auf der vollen Algebra. Die Familie aus \eqref{eq:biquaternion} ist im Allgemeinen nicht unter Multiplikation abgeschlossen. Beispielsweise ist \(e_1^2=-1\) und hat reellen Skalarteil. Die Identität benutzt \(\overline{q_1q_2}=\bar q_2\bar q_1\) und die Zentralität der Norm. Für die Konjugation \(q\mapsto uqu^{-1}\) mit \(N(u)=1\) bleibt der Skalarteil \(it\,1\) fest. Über die Matrixrealisierung der Biquaternionen ist \(N\) die Determinante und die Norm-eins-Gruppe \(SL(2,\mathbb C)\). Ihre Konjugation auf dem spurfreien Vektorteil ergibt die Abbildung nach \(SO(3,\mathbb C)\), mit Kern \(\{\pm 1\}\) \cite{feldkern}. Auf diesem Teil können wir damit die Lorentz-Feldinvariante aus \eqref{eq:field-invariant} im selben algebraischen Träger behalten.

Dabei treten zwei Nullbedingungen auf. Für das Feld gilt \(s(\mathbf z)=0\), für den gesamten Kern \(N(q)=0\), also \(s(\mathbf z)=t^2\). Auf einem lichtartigen Feld ist \(N(q)=-t^2\). Der Parameter \(t\) bleibt in dieser Feldrepräsentation ein Lorentzskalar. Eine physikalische Deutung des gesamten Kerns muss angeben, welche der beiden Nullbedingungen sie verwendet und welche Rolle der Skalar dabei spielt \cite{feldkern,grundidee}.

\paragraph{Ein lichtartiges Feld in Spinorkoordinaten.}
Für \(\psi=(\psi_1,\psi_2)\in\mathbb C^2\) parametrisiert der folgende Ausdruck die lichtartige Feldklasse.
\begin{equation}
 \mathbf z(\psi)=
 (\psi_1^2-\psi_2^2,\,i(\psi_1^2+\psi_2^2),\,-2\psi_1\psi_2).
 \label{eq:field-null-spinor}
\end{equation}
Die beiden ersten Quadrate ergeben zusammen \(-4\psi_1^2\psi_2^2\), das dritte gibt den Gegenbeitrag. Die Welle aus \eqref{eq:field-example} entspricht \(\psi=(\sqrt a,0)\) in den dort verwendeten Feldkoordinaten. Jedes von null verschiedene \(\mathbf z\) mit \(s(\mathbf z)=0\) hat die beiden Urbilder \(\pm\psi\). Man gewinnt sie aus
\[
 \psi_1^2=\tfrac12(z_1-iz_2),\qquad
 \psi_2^2=-\tfrac12(z_1+iz_2),\qquad
 \psi_1\psi_2=-\tfrac12 z_3 .
\]
Das ist eine algebraische Beschreibung der klassischen Feldklasse. Für ein quantisiertes Einphotonensystem wären Zustand und Observablen zusätzlich zu spezifizieren \cite{feldkern}.

So erhält der Ansatz einen konkreten algebraischen Kandidaten. Auf seinem Feldteil lassen sich die Feldklassen und die reell-lineare Austauschoperation beschreiben. Für die nächste Aufgabe brauchen wir einen Zustandsraum mit einem Skalarprodukt und einer erklärten Auswahl der Materiezustände. Der folgende Vergleich untersucht, wie eine solche Auswahl drei geeignete Zustände liefern könnte. Welche physikalische Darstellung sie realisiert, muss eigens gezeigt werden.

\subsection{Materiezustände und Generationen}
\label{sec:matter}

Ein Generationsindex soll Zustände unterscheiden, die in den relevanten inneren Ladungen übereinstimmen. Aus einem dreidimensionalen Unterraum allein folgt daher zunächst nur die Dimension. Ob die drei Richtungen Farben, Generationen oder eine andere Multiplizität bezeichnen, entscheidet die Wirkung der inneren Symmetrien und der Observablen auf diesen Zuständen.

Eine Vergleichskonstruktion des Companion macht die Wahl des Unterraums explizit \cite{companion,physics}. Sei \(T_0:\mathbb C^3\to\mathcal K\) eine Isometrie in einen komplexen Hilbertraum und \(P_c\) der orthogonale Projektor auf die im Vergleichsmodell zulässigen Zustände. Die Matrix \(A_0=T_0^*P_cT_0\) erfüllt
\[
 \langle v,A_0v\rangle=\|P_cT_0v\|^2 .
\]
Die Bedingung \(A_0>0\) besagt, dass die drei projizierten Ausgangszustände linear unabhängig bleiben. Dann ist
\begin{equation}
 S=P_cT_0A_0^{-1/2}, \qquad
 S^*S=A_0^{-1/2}T_0^*P_c^2T_0A_0^{-1/2}=\id_3,\qquad P_cS=S .
 \label{eq:seed}
\end{equation}
Damit erhält man drei orthonormale Zustände im Bild von \(P_c\). Die Wahl von \(T_0\) und \(P_c\) geht als Voraussetzung ein. \(P_c\) setzt die Zulässigkeit im Vergleichsraum um, \(C\) aus \eqref{eq:projections} die HC-Gradierung.

Ein Operator \(O\) auf \(\mathcal K\) lässt sich zu \(S^*OS\) komprimieren, sofern das Bild von \(S\) in seinem Definitionsbereich liegt. Erhält \(O\) dieses Bild, gilt \(OS=S(S^*OS)\). Die Kompression stellt dann die Einschränkung des Operators dar. Im allgemeinen Fall erfasst sie die Matrixelemente zwischen den drei ausgewählten Zuständen, während Übergänge aus ihrem Unterraum heraus zusätzlich zu untersuchen sind. Das betrifft insbesondere die Zuordnung eines Massenspektrums zu einem nativen Operator.

\paragraph{Eine positive Matrix aus der gewählten Zustandsauswahl.}
Der Companion untersucht eine Abbildung aus dem ausgewählten Zustandsraum in dessen orthogonales Komplement. Sei \(\widehat R\) ein beschränkter, dimensionslos normierter Operator auf \(\mathcal K\). Auf den Zuständen \(\psi_i=S\mathbf e_i\) setzt der Vergleich \(B_\perp=(I-P_c)\widehat R S\) an. Die daraus gebildete Matrix ist
\begin{equation}
 Q_\perp=B_\perp^*B_\perp,\qquad
 (Q_\perp)_{ij}
 =\big\langle (I-P_c)\widehat R\psi_i,(I-P_c)\widehat R\psi_j\big\rangle .
 \label{eq:generation-defect-readout}
\end{equation}
Für \(v\in\mathbb C^3\) ist \(\langle v,Q_\perp v\rangle=\|B_\perp v\|^2\). Außerdem gilt \([\widehat R,P_c]\psi_i=(I-P_c)\widehat R\psi_i\). Dessen Erwartungswert in \(\psi_i\) verschwindet wegen der Orthogonalität zum Bild von \(P_c\), auch wenn die Norm des Vektors positiv ist. Eine Auslese durch den Erwartungswert und eine Auslese durch die Übergangsnorm liefern damit unterschiedliche Größen. Bei exakter Kommutation \([\widehat R,P_c]=0\) verschwindet auch die gesamte Matrix \(Q_\perp\) \cite{companion,repo}.

Welche Zahlen würde man bei \(Q_\perp\ne0\) als Massenkandidaten lesen? Setzt man im Vergleich eine Skala \(\kappa>0\) und die effektive Matrix \(\kappa Q_\perp\) an, sind dafür ihre Eigenwerte relevant. Die Diagonale in den ursprünglichen Generationenkoordinaten enthält die drei quadrierten Normen der einzelnen Übergänge. Wie weit beides auseinanderliegen kann, zeigt die im Companion angegebene Rang-eins-Matrix \(Q=\mathbf1\mathbf1^*/3\) mit \(\mathbf1=(1,1,1)^{\mathsf T}\). Ihre Diagonale ist \((1/3,1/3,1/3)\). Aus \(Q\mathbf1=\mathbf1\) und \(Q|_{\mathbf1^\perp}=0\) folgen die Eigenwerte \((1,0,0)\). In der Koide-Kennzahl aus Abschnitt \ref{sec:mass} liefern diese beiden Auswertungen folglich \(1/3\) und \(1\). Der Vergleich macht sichtbar, welche Ausleseregel, welcher Operator und welche Skala feststehen müssen, bevor ein Massenspektrum bestimmt werden kann \cite{companion,repo}.

Eine physikalische Generationskonstruktion müsste aus dem gemeinsamen Träger den Ausgangsraum auswählen, die Ladungsdarstellung auf ihm bestimmen und zeigen, wie Massen die drei Zustände unterscheiden. Gleichung \eqref{eq:seed} legt dabei offen, an welcher Stelle die Wahl der Zustände eingeht. Der Koide-Test in Abschnitt \ref{sec:mass} setzt daran an, sobald ein Kandidat für die Massenauslese feststeht.

\subsection{Innere Symmetrien auf den Materiezuständen}
\label{sec:gauge}

Der bisherige Träger \(\mathbb C\otimes\mathbb H\) liefert zwei innere Wirkungen. Einheitsquaternionen liefern in einer quaternionischen Darstellung eine \(SU(2)\)-Wirkung. Komplexe Multiplikation liefert eine Phasenwirkung. Für eine Farbwirkung nimmt Dixon die Oktonionen als dritten Faktor hinzu und arbeitet auf \(\mathbb C\otimes\mathbb H\otimes\mathbb O\) \cite{dixon}. Hält man eine imaginäre Einheit \(e_7\) fest, ordnet sich der Oktonionraum relativ zu ihr in Singulett und Triplett. Furey gewinnt darauf die Ladungen \(0,\tfrac13,\tfrac23,1\) aus einem Zahloperator \cite{furey}. Der Companion übernimmt beides als Vergleich, nicht als Herleitung (Physik-Ausgabe, Teil „What a representation construction must specify“ \cite{physics}). Eine physikalische Eichzuordnung muss diese Wirkungen auf denselben Zuständen angeben, die im Massensektor verwendet werden.

Wann verträgt sich eine innere Symmetrie mit dem Kompensator? Dafür brauchen wir zuerst eine gemeinsame Darstellung. In Abschnitt \ref{sec:field-compensation} wirkte \(J_{\mathrm B}\) reell-linear auf dem Feldraum. Im Materiesektor arbeiten wir stattdessen mit einer reellen Gradierung. Ihre Involution \(J\) setzen wir komplex-linear auf die Komplexifizierung fort. Die Übertragung vom Feldraum in den Materiesektor ist eine eigene, deklarierte Abbildung.

Sei nun \(\rho(g)\) eine unitäre innere Symmetrie in dieser Darstellung. Kommutiert \(\rho(g)\) mit \(J\), so kommutiert sie mit beiden Projektoren aus \eqref{eq:projections}. Ihre Wirkung erhält dann die beiden graduierten Unterräume. Für die durch \(S\) ausgewählte Generationenschar aus \eqref{eq:seed} braucht man zusätzlich die Invarianz ihres Bildes unter \(\rho(g)\). Unter dieser Voraussetzung ist \(S^*\rho(g)S\) eine Darstellung auf dem ausgewählten Zustandsraum. Damit lassen sich die Ladungen direkt auf dem Raum prüfen, der auch die massenrelevanten Zustände enthält.

Eine gemeinsame komplexe Phase wirkt in der betreffenden Darstellung mit einem skalaren Generator. Für die elektroschwache Zuordnung braucht man die physikalischen Ladungsgewichte und die chiralen Darstellungen der Felder. Die Phasenwirkung legt diese Daten nicht fest. Ebenso muss eine Farbdarstellung über ihre Wirkung auf den Zuständen identifiziert werden.

\subsection{Wo ein Generationsindex sitzen könnte}
\label{sec:charged-multiplicity}

Eine plausible Ladungszuordnung stellt für die Massenauslese eine weitere Frage. Welche Operationen können zwischen möglichen Generationen unterscheiden und zugleich die vorgegebenen Ladungen erhalten? Dafür betrachten wir die Operatoren, die mit den deklarierten Lorentz- und Eichgeneratoren kommutieren. Ihre Wirkung innerhalb einer Multiplizität kann Zustände unterscheiden, die dieselbe Darstellung der Ladungen tragen.

Der Companion prüft diese Frage am Dixon-Vergleich \(\mathbb C\otimes\mathbb H\otimes\mathbb O\) \cite{dixon}, also an dem um die Oktonionen erweiterten Träger aus Abschnitt \ref{sec:gauge} (Runde v4.10.0-r19 \cite{dixon-commutant}). Er legt die Lorentzwirkung auf dem Quaternionfaktor, die schwachen Generatoren mit Chiralitätsprojektor, den Hyperladungsoperator und die Farbderivationen fest, die die Oktonionrichtung \(e_7\) erhalten. Auf einer Generation erhält man in diesem Vergleich die sechs Weyl-Multipletts \(Q,L,u^c,d^c,e^c,\nu^c\) mit komplexen Dimensionen \(12,4,6,6,2,2\). Der komplex-lineare Kommutant ist
\begin{equation}
 \operatorname{End}_{\mathrm{Lorentz},\mathrm{Eich}}(\mathcal K_{\rm Gen})
       \cong \mathbb C^6 .
 \label{eq:gauge-one-copy}
\end{equation}
In der deklarierten Darstellung sind die sechs irreduziblen Bestandteile paarweise inequivalent und treten mit Multiplizität eins auf. Nach dem Schur-Lemma wirkt ein kommutierender Operator auf jedem als Skalar. Die konkrete Generatorprüfung findet im reellen Kommutanten außerdem keinen von null verschiedenen antilinearen Anteil \cite{dixon-commutant}.

Diese sechs Skalare liefern Gewichte für verschiedene Multipletts. Ein Generationsindex braucht innerhalb einer Ladungsdarstellung mehrere Zustände. Für \(n\) identische Kopien hat man die Zerlegung
\[
 \mathcal K_{\rm Gen}^{\oplus n}
   \cong \bigoplus_{a=1}^{6}(\mathcal K_a\otimes\mathbb C^n).
\]
Auf dem \(a\)-ten Summanden ist die Wirkung \(\rho_a(g)\otimes I_n\). Damit kommutiert jeder Operator \(I_{\mathcal K_a}\otimes M_a\) mit \(M_a\in M_n(\mathbb C)\). Das Spektrum kann sich also innerhalb einer Ladungsmultiplizität ändern. Die Zahl \(n\) ist dabei gesetzt, nicht hergeleitet. Drei Generationen erhält man in diesem Vergleich nur, indem man \(n=3\) setzt \cite{dixon-commutant,dixon-copies}. Die Herkunft der Zahl drei bleibt Teil des offenen Problems OP3 zu Generationen und Koide-Relation (Abschnitt~\ref{sec:open-points}).

Was folgt daraus für die Feldoperationen aus Abschnitt \ref{sec:frames}? In der deklarierten Materiedarstellung wirkt die Multiplikation mit \(i\) zentral. Die komplexe Konjugation und der Austausch \(J_{\mathrm B}\) kommutieren dort mit Teilen der angesetzten Generatoren nicht. Ein ladungserhaltender Generationsoperator, der aus dem HC hervorgehen soll, braucht deshalb eine explizite Übertragung auf die Multiplizität. Seine Wirkung auf das elektromagnetische Dublett allein lässt diese Übertragung offen \cite{dixon-commutant}.

Kann die Farbe selbst den Generationsindex tragen? Eine Prüfung im Companion verneint das für lineare Auslesen. Legt man drei
farbtriviale Generationsrichtungen neben die komplexe Linie aus \(1,e_7\) und das Farbtriplett mit den Spalten \(e_4,e_5,e_6\), so wirkt
das nichttriviale Farbzentrum auf dem farbigen Teil als Faktor \(\omega=(-1+i\sqrt3)/2\). Eine farbäquivariante lineare Auslese \(\ell\) in ein
Singulett erfüllt dort \(\omega\ell=\ell\) und verschwindet deshalb, auch auf den drei Diagonalplätzen. Mit einem dualen Farbindex bleibt nur die
Spur als invariante Auslese, vom Rang eins. Die farbneutrale Restspalte hat dagegen Rang drei und trägt in der Lesart des Companion je
Zeile ein schwaches Dublett. Für lineare farbäquivariante Auslesen kann ein Generationsindex also nicht in der Farbe sitzen; er muss
dann in einer farbneutralen Multiplizität liegen, wie sie der Kopienvergleich oben ansetzt. Nichtlineare Auslesen erfasst die Prüfung
nicht. Die physikalische Zuordnung der Ladungen ist damit nicht hergeleitet, aber für diese Klasse von Auslesen ist ein Kandidat
ausgeschlossen (H113, Runde v4.10.0-r30 \cite{h113,companion}).

Damit ist die Prüfung der in \eqref{eq:seed} konstruierten Zustände konkreter. Neben ihrer Anzahl und Orthogonalität ist die Darstellung der Ladungen auf ihrem Bild anzugeben. Erst für den so spezifizierten Raum kann das Spektrum eines Massenoperators als Aussage über Generationen gelesen werden.

\subsection{Lokale Symmetrie und Kopplung}

Für \(U(x)=\rho(g(x))\) und \(D=d+A\) ergibt lokale Eichkovarianz \(D'(U\psi)=U D\psi\) die Anforderung
\[
 A'=UAU^{-1}-(dU)U^{-1}.
\]
Diese Transformationsregel legt fest, was ein vorgeschlagener Verbindungsoperator leisten muss. Seine konkrete Form, seine Felddynamik und seine Kopplungen an die ausgewählten Zustände sind weitere Bestandteile des Modells.

Für eine ungebrochene innere Symmetrie muss die Massenobservable auf dem ausgewählten Raum invariant unter der entsprechenden Ladungswirkung sein. Bei einer Kopplung an einen symmetriebrechenden Hintergrund sind die Transformation des Hintergrunds und die der Materiefelder gemeinsam zu berücksichtigen. So treffen sich Eich- und Massensektor in einer überprüfbaren Forderung an dieselben Zustände und Operatoren. Im recherchierten Stand bleiben die physikalischen Einbettungen, die Chiralität und die Kopplungsstärken offen.

\section{Was davon bei den Massen ankommt}
\label{sec:mass}

Die bisherige Konstruktion legt fest, wie man Zustände auswählt und unter welchen Bedingungen sie eine innere Symmetrie tragen. Das kann man schon für einen vorgegebenen Zustandsraum prüfen. Jetzt kommt eine andere Anforderung hinzu. Aus diesen Zuständen sollen konkrete Massen folgen. An den drei geladenen Leptonen können wir fragen, welche Eigenschaft eine gemeinsame Auslese haben müsste.

\subsection{Die Normbedingung im Raum der Massenwurzeln}

Für drei nichtnegative Massenwurzeln \(x_i=\sqrt{m_i}\), \(i=1,2,3\), mit \(x\ne0\) definieren wir
\begin{equation}
 K(x)=\frac{\sum_i x_i^2}{(\sum_i x_i)^2}.
 \label{eq:koide}
\end{equation}
Die physikalische Anwendung betrifft die drei geladenen Leptonen. Für die folgende Identität genügt ein beliebiger positiver Summenwert. Mit \(a=(\sum_i x_i)/3>0\) schreiben wir \(x=a(1,1,1)+d\), wobei \(\sum_i d_i=0\). Die beiden Anteile sind orthogonal. Folglich
\begin{equation}
 K(x)=\frac{3a^2+\|d\|^2}{9a^2}
      =\frac13+\frac{\|d\|^2}{9a^2},
 \qquad
 K(x)=\frac23\ \Longleftrightarrow\ \|d\|^2=3a^2 .
 \label{eq:equal-norm}
\end{equation}
Die Koide-Bedingung verlangt gleiche quadratische Normen des konstanten Anteils und der Abweichung. Ein zyklischer Basiswechsel erhält beide Normen und lässt ihr Verhältnis frei. Als Beispiel ergeben \(x=(1,2,3)\) den Wert \(7/18\) und drei gleiche Wurzeln den Wert \(1/3\). Die Rechnung verlegt die physikalische Frage damit auf die Auswahl der Normrelation \cite{physics}.

Eine für den Vergleich mit dem Spektrum nützliche Koordinate auf der
Nullsummenebene ist
\begin{equation}
 x_k=a\bigl[1+r\cos(\theta+2\pi k/3)\bigr],
 \qquad k=0,1,2,\qquad r\geq0 .
 \label{eq:koide-radius-phase}
\end{equation}
Die beiden Fourier-Vektoren mit Komponenten \(\cos(2\pi k/3)\)
und \(\sin(2\pi k/3)\) sind orthogonal und haben beide die
quadratische Norm \(3/2\). Damit gilt
\(\|d\|^2=\tfrac32a^2r^2\), und \eqref{eq:equal-norm} wird zu
\[
 K=\frac13+\frac{r^2}{6},\qquad K=\frac23\ \Longleftrightarrow\ r=\sqrt2 .
\]
Die Normrelation wählt einen Radius im Raum der normierten Wurzeln.
Entlang dieses Radius bleiben \(\theta\) und die gemeinsame
Normierung \(a\) als Modellgrößen zu bestimmen. Zusätzlich muss jeder
der drei Ausdrücke in eckigen Klammern nichtnegativ sein. Ist einer dieser Ausdrücke
negativ, liefert die Hauptwurzel der quadrierten Größe ihren
Betrag. Dann ändert sich der Nenner in \eqref{eq:koide}, und die
Radiusformel kann für dieses Spektrum nicht verwendet werden
\cite{companion,physics}.

\subsection{Der Test am Kompensator}

Auf dem euklidischen Wurzelraum sei \(e=(1,1,1)/\sqrt3\) und \(C_m=(\id+J_m)/2\) mit orthogonaler Involution \(J_m\). Angenommen, \(x=\alpha C_m e\) mit \(\alpha>0\) sei nichtnegativ und von null verschieden. Dann gilt \(\langle e,C_m e\rangle=\|C_m e\|^2\). Setzt man \(h=\langle e,J_m e\rangle\), folgt
\begin{equation}
 K(x)=\frac{1}{3\|C_m e\|^2}
      =\frac{2}{3(1+h)},\qquad
 K(x)=\frac23\ \Longleftrightarrow\ h=0.
 \label{eq:hc-criterion}
\end{equation}
Der Test verlangt also, dass \(e\) und \(J_me\) orthogonal sind. Die Projektoraxiome erlauben auch \(J_m=\id\), mit \(h=1\). Um die Orthogonalität physikalisch zu begründen, muss \(J_m\) aus dem gemeinsamen Träger folgen.

Dafür betrachtet die Physik-Ausgabe eine Koisometrie \(B:H\to\mathbb R^3\) mit \(BB^*=\id_3\) und \(J\)-invariantem Kern \cite{physics}. Da \(J\) selbstadjungiert ist, bleibt auch \((\ker B)^\perp=\operatorname{im} B^*\) invariant. \(B\) überträgt die Involution dieses orthogonalen Komplements als \(J_m=BJB^*\) auf den Wurzelraum, sodass \(BC=C_mB\). Ein unabhängig bestimmter Ausgangszustand \(f_0\) mit \(Bf_0=e\) und das physikalische Wurzelgesetz \(x=\alpha BCf_0\) würden die Bedingung auf
\[
 \langle B^*e,J B^*e\rangle=0
\]
zurückführen. Das benennt eine konkrete Prüfung am Träger. Die Auswahl von \(B\), \(f_0\) und des Wurzelgesetzes geht als eigene Voraussetzung ein. Ein rein ungerader Zustand liegt im Kern von \(C\) und liefert über diese direkte lineare Auslese keine von null verschiedenen Wurzeln.

Ein vollständig bestimmter Vergleich zeigt, welche Information in
der Wahl von \(J_m\) steckt. Wähle
\(v=(2,-1,-1)/\sqrt6\) und lasse \(J_m\) die orthonormalen
Vektoren \(e\) und \(v\) vertauschen. Auf ihrem orthogonalen
Komplement wirke \(J_m\) als Identität. Dann ergibt sich
\[
 C_m e=\frac12(e+v),\qquad
 x=\frac{\alpha}{2\sqrt3}
       \left(1+\sqrt2,\ 1-\frac1{\sqrt2},\ 1-\frac1{\sqrt2}\right).
\]
Das sind drei positive Hauptwurzeln im Vergleich, von denen zwei
gleich sind. Der Faktor vor der Klammer ist ihr
Mittelwert \(a\). Die erste Komponente hat die Richtung ausgewählt,
die in \eqref{eq:koide-radius-phase} der Phase \(\theta=0\)
entspricht. Hier ist also \(K=2/3\) bereits erreicht, bevor überhaupt
drei verschiedene Werte auftreten. Im Companion dient die
Konstruktion als Beispiel für die Realisierbarkeit von
\(\langle e,J_m e\rangle=0\). Für eine physikalische Erklärung
braucht man eine Regel, durch die der Träger gerade eine solche
Richtung und eine bestimmte Aufspaltung auswählt \cite{companion,physics}.

\subsection{Die Koide-Überlappung}

Die Überlappungsdarstellung setzt drei normierte Zustände \(v_i\) mit Gram-Matrix \(G_{ij}=\langle v_i,v_j\rangle\) an. In diesem Modell sind die Wurzeln \(x_i=s\lambda_i(G)\) mit \(s>0\). Da \(G\) positiv semidefinit ist und \(\Tr G=3\) gilt, ergibt sich
\[
 \Tr G^2=3+2\sum_{i<j}|\langle v_i,v_j\rangle|^2.
\]
Definieren wir \(\rho^2=\frac13\sum_{i<j}|\langle v_i,v_j\rangle|^2\), dann folgt
\begin{equation}
 K=\frac{\Tr G^2}{(\Tr G)^2}
   =\frac{1+2\rho^2}{3},
 \qquad K=\frac23\ \Longleftrightarrow\ \rho^2=\frac12 .
 \label{eq:overlap}
\end{equation}
Die Gleichung setzt das Wurzelgesetz für das Spektrum voraus. Sie bestimmt den Mittelwert der drei quadratischen Überlappungen. Paarweise gleiche Überlappungen sind eine zusätzliche Spezialisierung.


\subsubsection{Ein durchgehendes Beispiel mit veränderlicher Phase}

Wie kann bei festem \(K\) aus dem Beispiel mit einer doppelten
Wurzel ein Spektrum mit drei verschiedenen Wurzeln werden? Für die
im Companion untersuchte Vergleichsfamilie seien alle drei
paarweisen quadratischen Überlappungen gleich \(1/2\). Schreibe
\(U(z_1,z_2,z_3)=(z_2,z_3,z_1)\) und setze
\begin{equation}
 G(\varphi)=I_3+\frac1{\sqrt2}\left(
   e^{i\varphi/3}U+e^{-i\varphi/3}U^*\right).
 \label{eq:koide-circulant-gram}
\end{equation}
Das Produkt \(G_{12}G_{23}G_{31}=e^{i\varphi}/(2\sqrt2)\)
speichert die Bargmann-Phase \(\varphi\). Bei Umphasierung der
drei Zustände, \(v_j\mapsto e^{i\alpha_j}v_j\), wird
\(G_{ij}\) mit \(e^{i(\alpha_j-\alpha_i)}\) multipliziert.
Im Produkt heben sich diese Faktoren auf. Die Phase bleibt also
ein Merkmal des Tripels, auch wenn man die Phasen seiner drei
Vektorrepräsentanten frei wählt. Die zirkulante Darstellung
verteilt diese Phase gleichmäßig auf die drei gerichteten
Einträge \cite{physics,companion}.

Die Fourier-Basis diagonalisiert \(U\). Mit
\(\omega=e^{2\pi i/3}\) gehört zum Eigenwert \(\omega^k\) von
\(U\) die Eigenzahl
\begin{equation}
 \lambda_k(\varphi)=1+\sqrt2\cos\left(\frac{\varphi}{3}
                                     +\frac{2\pi k}{3}\right)
 \label{eq:koide-gram-spectrum}
\end{equation}
von \(G\). Für positiv semidefinites \(G\) sind dies
die in \eqref{eq:overlap} verwendeten Wurzeln in Einheiten von
\(s\). Ihr Mittelwert ist eins. Für \(\varphi=0\) ergibt sich
genau das normierte Wurzeltripel aus dem Projektorvergleich. In
der zirkulanten Familie entspricht die Änderung von \(\varphi\)
einer Änderung von \(\theta=\varphi/3\) in
\eqref{eq:koide-radius-phase}. Der Radius bleibt dabei \(\sqrt2\).
Somit kann man an diesem Beispiel direkt verfolgen, wie dieselbe
quadratische Überlappung unterschiedliche Wurzelspektren zulässt.

\subsubsection{Warum der Oktant am Rand liegt}

Für eine Realisierung durch drei Zustände muss \(G\) positiv
semidefinit sein. Die Determinantenentwicklung liefert
\begin{equation}
 \det G=1-\tfrac32+2\operatorname{Re}(G_{12}G_{23}G_{31})
       =\tfrac12(\sqrt2\cos\varphi-1).
 \label{eq:det}
\end{equation}
Alle \(2\times2\)-Hauptminoren sind gleich \(1/2>0\).
Im Hauptwertbereich der Phase ist \(G\) daher für
\(|\varphi|<\pi/4\) positiv definit. Bei \(|\varphi|=\pi/4\)
hat die Matrix Rang zwei, und für \(|\varphi|>\pi/4\) erhält
man keine Gram-Matrix von Zustandsvektoren. Die Determinante
zeigt damit, an welcher Stelle der formale Ausdruck für das
Spektrum in \eqref{eq:koide-gram-spectrum} seine Interpretation
als drei nichtnegative Hauptwurzeln verliert.

\begin{figure}[!htbp]
\centering
\begin{tikzpicture}[>=Latex,x=9.6cm,y=4.25cm,font=\small]
  \fill[qs blue!10] (-.25,0) rectangle (.25,.255);
  \draw[qs axis,->] (-.57,0)--(.58,0) node[right] {\(\varphi/\pi\)};
  \draw[qs axis,->] (0,-.39)--(0,.30) node[above] {\(\det G\)};
  \foreach \x/\lab in {-.5/{-\frac12},-.25/{-\frac14},.25/{\frac14},.5/{\frac12}} {
    \draw[qs axis] (\x,-.012)--(\x,.012);
    \node[below=2pt] at (\x,0) {\(\lab\)};
  }
  \draw[very thick,qs blue,domain=-.45:.45,samples=81]
     plot (\x,{(sqrt(2)*cos(180*\x)-1)/2});
  \draw[densely dashed,qs axis] (-.25,0)--(-.25,.255) (.25,0)--(.25,.255);
  \fill[qs blue] (-.25,0) circle (1.65pt) (.25,0) circle (1.65pt);
  \node[qs ink,fill=qs blue!10,inner sep=2pt] at (0,.09) {positiv definit};
  \node[qs ink,anchor=west] at (.42,-.18) {\(\det G<0\)};
\end{tikzpicture}
\caption{Determinante der Gram-Matrix bei festen paarweisen quadratischen Überlappungen \(1/2\). Im gezeigten Phasenintervall entsprechen die beiden Nullstellen dem Rangabfall. Im markierten Bereich zwischen ihnen sind alle drei Eigenwerte positiv. In der gesamten zulässigen Familie hat die Koide-Auslese denselben Wert.}
\label{fig:koide-phase}
\end{figure}


Bei \(\varphi=\pi/4\) erhält man
\[
 \left(\lambda_0,\lambda_1,\lambda_2\right)
  =\left(\frac{3+\sqrt3}{2},\,0,\,\frac{3-\sqrt3}{2}\right).
\]
Den Rand kann man direkt in \(\mathbb C^2\) realisieren. Wähle
\(v_1=(1,0)\), \(v_2=(1,1)/\sqrt2\) und
\(v_3=(1,i)/\sqrt2\). Die drei Überlappungen haben den
Betrag \(1/\sqrt2\) und ihr gerichtetes Produkt ist
\((1+i)/4\). Es sind die drei Spin-\(\tfrac12\)-Zustände
eines Bloch-Oktanten, den die Physik-Ausgabe diskutiert.
Die positive Matrix im Inneren des Phasenintervalls hat
hingegen Rang drei und verlangt eine Realisierung in
mindestens drei komplexen Dimensionen. Eine mögliche
Realisierung liefern dann die drei Spalten von \(G^{1/2}\).
Für \(0<|\varphi|<\pi/4\) sind alle drei Wurzeln positiv
und verschieden. Der Wert in \eqref{eq:overlap} ist an
Rand und Innerem derselbe. Die Rechnung legt daher die
zulässige Phasenfamilie offen, während die Physik erklären
müsste, welcher Wert in dieser Familie tatsächlich
ausgewählt wird. Die Interpretation des Defizits
\(\pi/4-|\varphi|\) als Kern in Abschnitt~5.8 der
Physik-Ausgabe bleibt eine bedingte Konstruktion
\cite{physics,companion}.

\subsubsection{Welches Vorzeichen der Oktantschritt hat}

Der Oktant legt den Betrag der Phase fest, nicht ihr Vorzeichen. Durchläuft
man dieselben drei Zustände in umgekehrter Reihenfolge,
\(v_1\to v_3\to v_2\to v_1\), bleiben Anfangs- und Endstrahl und alle drei
Übergangswahrscheinlichkeiten \(\tfrac12\) gleich. Das gerichtete Produkt
wird aber zu \((1-i)/4\), der Schritt also zu \(-\pi/4\) statt \(+\pi/4\).
Beträge und Endpunkte entscheiden das Vorzeichen daher nicht. Auch eine
gleichmäßige Verteilung der aufsummierten Phasen sagt nichts: Die beiden
konstanten Schrittfolgen \(\pm1\) modulo acht durchlaufen in einem Zyklus
alle acht Phasen \cite{h116}.

Der Kompensator liefert das Vorzeichen ebenfalls nicht. Auf den
geschlossenen ungeraden Generationszuständen gilt \(JS=-S\), also
\(CS=(S+JS)/2=0\). Eine reelle Projektion vom Rang fünf, die eine reelle
Richtung der komplexen Nullgeraden der Oktant-Auslese entfernt, gibt es
zwar. Die beiden reellen Richtungen \(n\) und \(in\) dieser Geraden liefern
aber dieselben Ränge und dieselbe Auslese, und \(J\) wählt nicht zwischen
ihnen \cite{h115}.

Übrig bleibt die Auswahl des geordneten Rückwegs. Ein Modell mit zwei
Fenstern macht den Zeitpunkt sichtbar. Das erste Fenster hält die
zulässigen Orientierungen als Maske in einem vorläufigen Zustand, das
zweite prüft sie mit einer zweiten Maske, und ein Vorzeichen entsteht erst,
wenn genau eine Orientierung übrig bleibt. Von den sechzehn Maskenpaaren
ergeben drei \(+\pi/4\), drei \(-\pi/4\), neun sind gesperrt und eines
bleibt mehrdeutig. Die vorhandenen Regeln des Companion vertragen sich mit
diesem Zeitpunkt, enthalten aber keine Größe, die zwischen den beiden
Vorzeichen entscheidet \cite{h117}. Alle drei Rechnungen waren vorab
festgelegt, ihr Ausgang stand aber analytisch fest; sie gelten deshalb
nicht als blinde Tests. Die Physik-Ausgabe von v5.3.0 fasst sie in Abschnitt~5.9 zusammen \cite{physics}.
Eine spätere Prüfung ordnet die Orientierung ein: Frame-Abbildungen und
Parität erhalten sie, die Zeitumkehr der Feldamplituden wirkt als komplexe
Konjugation und kehrt sie um. Das Vorzeichen ist also eine Zeitorientierung
der Auslese, und keine im Companion festgehaltene Struktur legt es fest
\cite{h118}.

\subsubsection{Was sich beim Kombinieren von Präparationen ändert}

Die Phase kann nur dann als Eigenschaft einer Auslese
verwendet werden, wenn feststeht, an welcher Stelle der
Auslese sie eingeht. Betrachten wir zwei der oben
beschriebenen positiven Tripel mit \(\theta_\pm=\theta_0\pm h\)
im selben positiv zulässigen Intervall, demselben \(a\)
und derselben Zuordnung der Komponenten in einem
festen Materieframe. Wenn die beiden Wurzeln einer
Komponente zuerst zu gleichen Teilen gemittelt werden,
folgt mit der Addition der beiden Kosinuswerte
\begin{equation}
 \bar x_k=a\bigl[1+\sqrt2\cos h\,
                      \cos(\theta_0+2\pi k/3)\bigr],\qquad
 K(\bar x)=\frac{1+\cos^2 h}{3}.
 \label{eq:koide-mixture}
\end{equation}
Die Präparationen dürfen also einzeln \(K=2/3\) erfüllen,
während die gemittelte Wurzelauslese für eine echte
Phasendifferenz einen kleineren Radius hat. Würde man die
beiden Massen \(x_{+k}^2,x_{-k}^2\) unmittelbar mitteln,
ergäbe sich pro Komponente
\[
 \frac{x_{+k}^2+x_{-k}^2}{2}
   =\bar x_k^{\,2}+\frac{(x_{+k}-x_{-k})^2}{4}.
\]
Das zweite Ergebnis enthält einen Varianzterm. Beide
Verfahren gehen von denselben Präparationen aus und
unterscheiden sich in ihrer Operationsreihenfolge.

Noch einmal eine eigene Möglichkeit ist die kohärente
Addition zweier Pfadbeiträge \(B_1+e^{i\phi}B_2\) zu
einem Übergangsoperator \(B\). Im Produkt \(B^*B\)
entstehen dann die Kreuzterme
\(e^{i\phi}B_1^*B_2+e^{-i\phi}B_2^*B_1\).
Eine relative Phase kann deshalb im Spektrum des
Operators wirken. Ein gemeinsamer Phasenfaktor vor
\(B\) fällt aus \(B^*B\) heraus. Für eine native
Zustandspräparation müsste man angeben, welche dieser
Operationen physikalisch realisiert wird. Die
Phasenmittelung ist im Companion ein ausdrücklich
bedingter Vergleich \cite{companion,physics}.

\subsection{Eine Vergleichsdynamik mit leichter und schwerer Komponente}

Die Kompression aus Abschnitt \ref{sec:matter} kann mit einem Übergangsoperator zu einem zusätzlichen Raum verbunden werden. Als konkreten Vergleich betrachten wir einen Übergang \(B_{\mathrm{lh}}\) zwischen zwei endlichdimensionalen Hilberträumen mit ihren deklarierten positiven Metriken. Der schwere Block habe die konstante Lücke \(\Delta>0\) und die Kopplung das Gewicht \(\mu>0\). In orthonormalen Basen lautet der hermitesche Blockoperator
\[
 \Omega=
 \begin{pmatrix}
 0 & \mu B_{\mathrm{lh}}^* \\
 \mu B_{\mathrm{lh}} & \Delta\id
 \end{pmatrix}.
\]
Die im Companion behandelte Form mit gewichteten Metriken geht durch einen Metrikwechsel in diese Form über \cite{companion,repo}. Gehört \(\sigma_i>0\) zu den Singulärwerten des Übergangs, reduziert sich der zugehörige Teil auf die Matrix mit Diagonale \(0,\Delta\) und Nebendiagonale \(\mu\sigma_i\). Die negative Eigenzahl \(-m_i\) liefert
\begin{equation}
 m_i(m_i+\Delta)=\mu^2\sigma_i^2 .
 \label{eq:gap}
\end{equation}
Für den Vergleich sind \(\sigma_i\) dimensionslos und \(\mu,\Delta,m_i\) haben die Dimension einer Masse. Für \(\mu\sigma_i\ll\Delta\) ist \(m_i\) in führender Ordnung \(\mu^2\sigma_i^2/\Delta\). Diese Näherung hängt von der Lücke ab. Der leichte Eigenzustand hat bei von null verschiedener Kopplung auch eine Komponente im schweren Block.

Das zugehörige Eigenwertproblem lässt sich mit nur zwei
Komponenten verfolgen. Seien \(u_i,w_i\) normierte
Singularvektoren mit \(B_{\mathrm{lh}}u_i=\sigma_iw_i\)
und \(B_{\mathrm{lh}}^*w_i=\sigma_iu_i\). Für den Eigenzustand
zur Eigenzahl \(-m_i\) sind die Koeffizienten in dieser Basis
proportional zu
\[
 \left(1,-\frac{\mu\sigma_i}{\Delta+m_i}\right).
\]
Damit ist der quadrierte Anteil der schweren Komponente
\(m_i/(\Delta+2m_i)\), nach Einsetzen von \eqref{eq:gap}
in die Normierung. Ein frei gewählter Übergang liefere etwa
\(q_i=\sigma_i^2\). Mit der dimensionslosen Kopplung
\(\eta=\mu/\Delta\) folgt exakt
\begin{equation}
 m_i=\frac{\Delta}{2}\left(\sqrt{1+4\eta^2q_i}-1\right).
 \label{eq:koide-scale-gap}
\end{equation}
Bei festem \(q_i,\eta\) multipliziert eine Änderung von
\(\Delta\) alle \(m_i\) mit demselben Faktor. Ändert man
\(\eta\) bei festem \(q_i\), ändern sich im Allgemeinen auch
die Massenverhältnisse. Der Parameter für
die relative Wirkung der Kopplung und die absolute
Massenskala haben damit getrennte Aufgaben. Der Anteil
der schweren Komponente im Eigenzustand hängt von der
dimensionslosen Wahl ab und bleibt bei einer gemeinsamen
Skalierung von Lücke und Kopplung gleich
\cite{companion,repo}.

Die rechte Seite von \eqref{eq:gap} setzt den Übergangsoperator voraus. Die Gram-Auslese in \eqref{eq:overlap} setzt dagegen ein Gesetz für die \emph{Wurzeln} der Massen voraus. Eine Identifikation der Ergebnisse würde deshalb eine Abbildung benötigen, die aus dem Übergang und seinen Skalen genau das Wurzelspektrum erzeugt. So wird die native Operatoridentifikation zu einer konkreten Anforderung an zwei explizite Modelle.

\subsection{Was ein Massengesetz zusätzlich festlegen muss}

Im Gram-Beispiel stehen damit drei unterschiedliche Entscheidungen
an. Die Wahl der Überlappungen fixiert die Normrelation. Die Phase
bestimmt innerhalb des gezeigten Intervalls, welches Verhältnis
der Wurzeln erreicht wird. Der Faktor \(s\) gibt den Wurzeln ihre
Einheit und setzt die absolute Massenskala über \(s^2\). Der
Projektorvergleich liefert ein Kriterium für die Normrelation.
Die Blockdynamik liefert ein Beispiel dafür, wie in einer
anderen Auslese ein dimensionaler Parameter eingeht. Um diese
Stücke zu verbinden, müsste der gemeinsame Träger die
Zustandsauswahl aus Abschnitt \ref{sec:matter} und einen
nativen Operator liefern, dessen Wurzelspektrum gerade diese
Eigenschaften besitzt. Die Auswahl muss für einen Test vor dem
Vergleich mit den Zielmassen feststehen.

Im folgenden Abschnitt geht es um die nächste Stufe dieser
Frage. Ein vorgegebener Massenterm lässt sich in eine
dynamische Gleichung einsetzen. Wie der native Massenoperator
diesen Term liefern und wie derselbe Materiezustand in einer
geometrischen Quelle auftreten würde, wird damit zu einer
konkreten Anforderung an die Kopplung.

\section{Wie Materie zur Quelle der Geometrie werden könnte}
\label{sec:geometry}

Angenommen, das Modell könnte den ausgewählten Materiezuständen Massen zuordnen. Wir wissen dann noch nicht, wie dieselben Zustände auf die Geometrie wirken. Dazu braucht es eine Quelle und eine Gleichung für ihre Wirkung. Der erste Test ist, ob die ausgelesenen Größen eine geschlossene Dynamik haben. Erst dann stellt sich die Frage, welche Feldgleichung diese Dynamik beschreiben kann.

\subsection{Geschlossene Dynamik ausgelesener Variablen}

Die Zustandszerlegung aus Abschnitt \ref{sec:carrier} lässt sich auf eine zeitabhängige Auslese erweitern. Sei \(u\) ein Zustand eines endlichdimensionalen reellen Vergleichssystems mit \(\dot u=G(t)u\). Die Auslese sei \(r=T(t)u\) mit differenzierbarem surjektivem \(T(t)\). Zu einer Rechtsinversen \(R(t)\) von \(T(t)\) lässt sich jeder Zustand als \(u=Rr+z\) mit \(Tz=0\) schreiben. Dann gilt
\begin{align}
 \dot r&=(\dot T+TG)Rr+(\dot T+TG)z,\\
 \dot r\ \text{ist durch }(t,r)\text{ bestimmt}
      &\quad\Longleftrightarrow\quad (\dot T+TG)\ker T=0 .
 \label{eq:autonomy}
\end{align}
Die Notwendigkeit folgt, indem zu einem Zustand ein beliebiges Element des Kerns addiert wird. Die Hinreichendheit folgt aus dem Verschwinden des zweiten Terms. Der erste Term ist dann unabhängig von der Wahl der Rechtsinversen. Für ein konstantes \(T\) verlangt \eqref{eq:autonomy}, dass \(G\) den Kern erhält. Für eine zeitabhängige Auslese gehört ihre eigene Änderung zur Bilanz \cite{physics}.

Das Resultat betrifft einen linearen endlichen Vergleich. In einer Feldtheorie kommen unter anderem Definitionsbereiche von Differentialoperatoren, Randbedingungen, Eichfreiheit und die Fortpflanzung von Zwangsbedingungen hinzu. \eqref{eq:autonomy} macht die lokale Frage nach den ausgelassenen Zuständen sichtbar. Es entscheidet diese weiteren Fragen nicht.

\paragraph{Was ein verschwundener Anteil noch bewirken kann.}
Für die feste HC-Projektion aus \eqref{eq:projections} ist \(r=Cu\). Dann führt \eqref{eq:autonomy} auf \(CGP_-=0\). Als Beispiel vertausche \(J\) im Zustandsraum \(\mathbb R^2\) die beiden Koordinaten und sei \(G=\operatorname{diag}(1,0)\). Der Zustand \(u=(1,-1)^{\mathsf T}\) hat \(Cu=0\) und dennoch \(CGu=(1,1)^{\mathsf T}/2\). Aus dem ausgelesenen Zustand allein kann man also in diesem Vergleich seine Ableitung nicht bestimmen \cite{physics}.

\subsection{Coframe und Quelle in einem Cartan-Vergleich}

Eine mögliche geometrische Auslese ist ein reeller, punktweise invertierbarer Coframe \(E^a{}_\mu\). Dann lässt sich
\begin{equation}
 g_{\mu\nu}=\eta_{ab}E^a{}_\mu E^b{}_\nu,\qquad
 \eta=\operatorname{diag}(-1,1,1,1)
 \label{eq:metric}
\end{equation}
als Lorentzmetrik definieren. Eine bilineare Auslese für den Coframe muss daher vollen Rang haben.

Der explizite Einstein--Cartan-Vergleich des Companion setzt Coframe, Lorentzverbindung und ein Spinortripel als unabhängige Felder an. Ein Multiplikator koppelt den Coframe an eine bilineare Spinorauslese \(\mathcal C^a{}_\mu(\Psi,D\Psi)\). Die vorgeschriebene Wirkung hat die Form
\begin{equation}
 S_{\mathrm{vgl}}=\int\! d^4x\,
 \left[\left|\det E\right|\left(\tfrac c2 R(E,\omega)-u_0\right)
 +\mathcal L_{\mathrm D}
 +\lambda_a{}^\mu\bigl(E^a{}_\mu-\mathcal C^a{}_\mu\bigr)\right].
 \label{eq:cartan}
\end{equation}
Hier ist \(\mathcal L_{\mathrm D}\) die dort definierte Spinor-Lagrangedichte, \(\omega\) die unabhängige Lorentzverbindung, \(c\) ein Vergleichskoeffizient und \(u_0\) ein konstanter Term. Die Variation nach \(\lambda_a{}^\mu\) liefert \(E^a{}_\mu=\mathcal C^a{}_\mu\). Die Variationen nach Spinoren, Verbindung und Coframe stellen weitere Gleichungen. Insbesondere muss eine algebraisch bestimmte Verbindung mit ihren tatsächlich daraus folgenden Ableitungen eingesetzt werden, wenn man eine reduzierte Feldformulierung prüft \cite{companion,repo}.

In der konkreten neutralen Spinorrechnung wirken die beiden Vergleichsmatrizen \(Q\) und \(M_{\mathrm D}\) auf die drei Kanäle. Mit \(v=|\det E|\), \(e=E^{-1}\), \(K^a=\beta\Gamma^a\) und \(D_\mu\Psi=\partial_\mu\Psi-\tfrac14\omega_{\mu,ab}\Gamma^a\Gamma^b\Psi\) werden die Auslese und die Dichte durch
\begin{align}
 \mathcal C^a{}_\mu&=\operatorname{Im}(\Psi^\dagger QK^aD_\mu\Psi),\\
 \mathcal L_{\mathrm D}
 &=-v\operatorname{Im}(\Psi^\dagger e_a{}^\mu K^aD_\mu\Psi)
   -v\operatorname{Re}(\Psi^\dagger(M_{\mathrm D}\otimes\beta)\Psi)
\end{align}
festgelegt. Die Matrizen \(\Gamma^a,\beta\) bezeichnen die Dirac-Matrizen in der Vergleichskonvention. Kanal- und Spinindizes werden getrennt kontrahiert. Diese Festlegung macht sichtbar, welche Spinorgröße für die geometrische Auslese eingesetzt wird und an welcher Stelle ein Massenoperator in die Spinordynamik eingeht.

Die Vergleichswirkung setzt ihre Krümmung, Materiedichte und Kopplungskoeffizienten als Modellbestandteile an. Sie zeigt, wo eine Mate\-rie\-größe und der Coframe in derselben Variationsrechnung zusammentreffen. Für den nativen Übergang aus Abschnitt \ref{sec:mass} wäre zusätzlich zu bestimmen, welche dieser Eingaben aus dem gemeinsamen Träger folgen und wie die physikalische Quelle an die dortige Massenauslese gekoppelt ist.

\subsection{Von Daten an einem Punkt zu einer Feldlösung}

Die Cartan-Rechnungen im Companion prüfen endliche Ableitungsdaten an einem Punkt. Nach Eliminierung höherer freier Ableitungen ergeben sich lineare Bedingungen an vorgeschriebene Daten und den Quellenzuwachs. Ist die reduzierte Quellenmatrix \(L\) von vollem Spaltenrang, besitzt \(Lq=b\) für kompatible Daten genau eine Lösung \(q=L^+b\), wobei \(L^+\) die Moore--Penrose-Pseudoinverse ist. Die Kompatibilitätsbedingung lautet \((I-LL^+)b=0\). Diese Rechnung kann angeben, welche Daten an diesem Punkt eine gemeinsame Fortsetzung der geprüften endlichen Ordnung zulassen.

Der Vergleich enthält bereits eine bedingte Aussage zur Fortpflanzung der gravitativen Zwangsbedingungen. Sind die Verbindung kompatibel eliminiert und sowohl die Konstitutivgleichung \(E^a{}_\mu=\mathcal C^a{}_\mu\) aus der Variation nach \(\lambda_a{}^\mu\) als auch die Sondengleichung, die Euler-Gleichung der Spinoren, in einem offenen Gebiet erfüllt, geben die Ward-Identitäten dort einen symmetrischen und divergenzfreien Coframe-Resttensor. Mit den räumlich projizierten spurmodifizierten Coframe-Gleichungen erhält man unter den angegebenen Regularitäts- und Randbedingungen ein symmetrisch hyperbolisches System für die vier gravitativen Zwangsgrößen. Bei verschwindenden Anfangsdaten und ohne einlaufenden Zwangsfluss bleiben diese Größen null. Die Anwendung setzt die beiden zuerst genannten Feldgleichungen auf dem Gebiet voraus. Ein Verschwinden ihrer Reste an einem einzelnen Punkt reicht dafür nicht aus \cite{companion,repo}.

Eine Lösung in einem offenen Gebiet erfordert zusätzlich die gleichzeitige Erfüllung der Konstitutiv- und Sondengleichungen, kompatible Flächendaten und eine Kontrolle der Fortpflanzung der Zwangsbedingungen. Das bleibt eine eigene Aufgabe. Entsprechend unterscheidet der Companion die punktweise Prüfung des Cartan-Vergleichs von der allgemeinen Schließung des Feldsystems.

\subsection{Vom Materieantrieb zur geometrischen Antwort}
\label{sec:op5-coupling}

Im Cartan-Vergleich tritt ein vorgegebener Massenterm in der Spinor-Lagrangedichte auf. Für die Vereinheitlichungsfrage wollen wir weitergehen. Wie kann eine Änderung der Materie eine geometrische Antwort erzeugen, und woran lässt sie sich von der gewöhnlichen Gravitation des Aufbaus unterscheiden? Dazu verfolgen wir ein und dieselbe Modulation von der Materiequelle bis zur Auslese. Die Rechnung spielt im Materieframe \(M\) und benutzt Einheiten mit Lichtgeschwindigkeit eins. Als Eingabe dienen die Massenparameter des Dirac--Yukawa-Vergleichs. Die Identifikation dieser Parameter mit der nativen Massenauslese aus Abschnitt \ref{sec:mass} ist weiterhin offen.

Die beiden Wege in der folgenden Skizze setzen bei derselben Materialvariation an. Über den Energie-Impuls-Tensor ist ihr gewöhnlicher Beitrag bestimmt, sobald die mechanische Quelle und die Randdaten festliegen. Für den gerichteten Beitrag kommt zusätzlich die Zuordnung der radialen Higgsquelle zur Quelle des Vergleichsfelds hinzu. Genau diese Zuordnung wird unten als offene Annahme sichtbar.

\begin{figure}[!htbp]
\centering
\begin{tikzpicture}[>=Latex,font=\small,
  box/.style={draw=qs axis,rounded corners=2pt,align=center,text width=51mm,
             minimum height=10mm,inner sep=4pt,fill=white}]
  \node[box,text width=65mm] (density) at (0,0) {Präparierte Materialverschiebung\\\(\delta n_i\) und \(\delta T_{\mu\nu}\)};
  \node[box,fill=qs blue!7] (gravity) at (-3.55,-1.61)
     {Gewöhnliche Gravitation\\\(D_0\bar\gamma_{\mu\nu}=16\pi G\,\delta T_{\mu\nu}\)};
  \node[box,fill=qs rust!9] (higgs) at (3.55,-1.61)
     {Radiale Higgsquelle\\\(\delta\mathcal J_h\) aus W2};
  \node[box,fill=qs blue!7] (exterior) at (-3.55,-4.13)
     {Kompakter linearer\\mechanischer Beitrag};
  \node[box,fill=qs rust!9] (profile) at (3.55,-4.13)
     {Gerichtetes OP5-Profil\\\(\delta\mathcal J_{\rm OP5}=\mathcal P_{\rm dyn}\delta H\)};
  \node[box,fill=qs rust!9] (axis) at (3.55,-5.66)
     {Gezeitenkontrast auf dem Plateau\\\(\Delta\mathcal I\)};
  \draw[->,thick,qs blue] (density.south west) -- (gravity.north);
  \draw[->,thick,qs rust] (density.south east) -- (higgs.north);
  \draw[->,thick,qs blue] (gravity) -- (exterior);
  \draw[->,thick,dashed,qs rust] (higgs) -- (profile);
  \node[anchor=east,align=right,inner sep=4pt] at (3.52,-2.9)
      {angesetzte Quellenidentifikation\\[-1pt]\(\delta\mathcal J_{\rm OP5}=\delta\mathcal J_h\)};
  \draw[->,thick,qs rust] (profile) -- (axis);
\end{tikzpicture}
\caption{Die beiden Antwortwege derselben linearen Materialpräparation. Der gewöhnliche Beitrag folgt aus dem mechanischen Tensor. Der gerichtete Vergleich setzt die Kopplung der radialen W2-Quelle an das OP5-Profil voraus. Die gestrichelte Verbindung markiert diese Annahme.}
\label{fig:op5-modell}
\end{figure}


\subsubsection{Welche Quelle die Massenterme liefern}

W2 bezeichnet die gemeinsame Vergleichswirkung des Companion: Gravitation nach Einstein und Cartan mit den Spinoren einer Generation und ihren Eichfeldern (W1), ergänzt um ein Higgs-Dublett mit Yukawa-Kopplungen (Runden v4.9.0-r81 und v4.9.0-r82 \cite{companion}). Sei \(h_{\rm W}=v+\eta_h\) der Higgsradius von W2 bei festen Higgswinkeln und \(v>0\) sein Referenzwert. In der gewählten Normierung gilt \(m_i(h_{\rm W})=y_i h_{\rm W}\). Für nichtrelativistische Besetzungen mit Teilchenzahldichten \(n_i\) behalten wir zunächst die Ruheenergie. Dabei setzen wir die Ruheenergie des Materials mit seiner Yukawa-Masse gleich, nehmen also an, dass sie ganz vom Higgsradius getragen wird. Das ist eine Modellannahme dieser Fallstudie; Runde v4.10.0-r35 \cite{layers} führt den Higgsanteil der Ruhemasse als eigene Größe. Die Materiedichte und ihre radiale Quelle sind dann
\begin{align}
 \epsilon_{\rm mat}(h_{\rm W},n_i)&=h_{\rm W}\sum_i y_i n_i,
 &
 \mathcal J_h&=-\partial_{h_{\rm W}}\epsilon_{\rm mat}
              =-\sum_i y_i n_i .
 \label{eq:op5-yukawa}
\end{align}
Die Dichte \(n_i\) zählt bei beidseitiger Besetzung Teilchen und Antiteilchen gemeinsam. Die dem Higgsradius zugeordnete Quelle hängt dabei von der Besetzungssumme ab. Eine kleine Verschiebung \(\mathbf u\) einer homogenen Referenz ändert diese durch \(\delta n_i=-n_{i0}\nabla\cdot\mathbf u\). Deshalb ist
\begin{equation}
 \delta\mathcal J_h=\left(\sum_i y_i n_{i0}\right)\nabla\cdot\mathbf u,
 \qquad
 \left.\delta\epsilon_{\rm mat}\right|_v=-v\,\delta\mathcal J_h .
 \label{eq:op5-density}
\end{equation}
Kompression und Verdünnung liefern Quellenstörungen mit entgegengesetztem Vorzeichen. Sie ändern zugleich die gewöhnliche Energiedichte. Bewegt sich auch der Higgsradius, kommt zu \(\left.\delta\epsilon_{\rm mat}\right|_v\) der Beitrag \(\eta_h\sum_i y_i n_{i0}\) hinzu. Der folgende mechanische Abgleich verfolgt den direkten Transport der Ruhemasse bei festem Referenzradius. Die Feldenergie und der Impulsaustausch mit einem veränderlichen Higgsfeld gehören zur vollständigen Bilanz \cite{op5-density}.

Die Quelle in \eqref{eq:op5-density} beschreibt zunächst einen Dichteantrieb. Das lässt sich auch am Druck prüfen. Bei fester Entropie pro Teilchen gilt für eine einzelne Art \(p_{\rm mat}=n\partial_n\epsilon_{\rm mat}-\epsilon_{\rm mat}\). Die reine Ruheenergie ist proportional zu \(n\), sodass ihr Beitrag zu \(p_{\rm mat}\) verschwindet. Elastische und kinetische Terme sind deshalb als eigene Beiträge zu führen. Der lokale Energieabgleich bei festem \(v\) ersetzt die Kraftbilanz eines bewegten Materials nicht. Mit der Signatur \((-+++)\) und der Konvention \(T_{\mu\nu}=-2\,\delta S_{\rm mat}/(\sqrt{-g}\,\delta g^{\mu\nu})\) ergibt die Diffeomorphismusidentität auf den Materiegleichungen
\begin{equation}
 \nabla_\mu T^\mu{}_\nu=\mathcal J_h\,\partial_\nu h_{\rm W}
 \label{eq:op5-matter-ward}
\end{equation}
für einen als extern behandelten Higgsradius und ohne weitere äußere Kräfte. Bei dynamischem \(h_{\rm W}\) ist diese rechte Seite Teil eines Austauschs, den die übrigen Felder in der Gesamtbilanz aufnehmen müssen. Im mechanischen Vergleich unten werden der Referenzradius festgehalten und die inneren elastischen Kräfte gemeinsam bilanziert. Damit lässt sich \eqref{eq:op5-matter-ward} von der Erhaltung des präparierten linearen Tensors abgrenzen \cite{op5-density}.

\subsubsection{Wie sich der gewöhnliche Außenbeitrag kontrollieren lässt}

Kann man diese Quelle schalten und ihren gewöhnlichen gravitativen Beitrag zugleich kontrollieren? Der Vergleich benutzt zwei überlagerte isotrope elastische Kontinua mit Ruhedichten \(\rho_1,\rho_2>0\). Der lineare Cauchy-Tensor der ersten Komponente ist \(\sigma(\mathbf u)=\mathsf C\,\varepsilon(\mathbf u)\). Die zweite Komponente hat die Steifigkeit \(c_\sigma\mathsf C\) mit \(c_\sigma>0\). Wir schreiben \(\mathcal L\mathbf u=\nabla\cdot\sigma(\mathbf u)\) und
\begin{equation}
 A_1=\rho_1\partial_t^2-\mathcal L,\qquad
 A_2=\rho_2\partial_t^2-c_\sigma\mathcal L .
 \label{eq:op5-elastic}
\end{equation}
Bei konstanten Koeffizienten kommutieren \(A_1\) und \(A_2\). Für ein glattes, räumlich kompaktes Vektorfeld \(\mathbf b\) wähle
\[
 \mathbf u=A_2\mathbf b,\quad
 \mathbf w=-A_1\mathbf b,\quad
 \mathbf d=\mathbf u-\mathbf w+k^{-1}A_1\mathbf u .
\]
Eine innere Paarfeder mit \(k>0\) und vorgeschriebener relativer Ruhelage \(\mathbf d\) liefert die entgegengesetzten Kräfte \(\mp k(\mathbf u-\mathbf w-\mathbf d)\). Somit gelten während der Schaltung beide Bewegungsgleichungen \(A_1\mathbf u=-k(\mathbf u-\mathbf w-\mathbf d)\) und \(A_2\mathbf w=+k(\mathbf u-\mathbf w-\mathbf d)\). Die gezielte Änderung von \(\mathbf d\) ist eine Eingabe des Vergleichs \cite{op5-motion}.

Die beiden Gleichungen für die Feder lassen sich direkt an der gewählten Verschiebung kontrollieren. Es gilt \(A_1\mathbf u=A_1A_2\mathbf b\) und \(A_2\mathbf w=-A_2A_1\mathbf b=-A_1\mathbf u\). Beim Einsetzen von \(\mathbf d\) wird die Federkraft auf die erste Komponente mit \(A_1\mathbf u\) identifiziert. Damit ist die Summe ihrer Kräfte im Referenzmodell zu jedem Zeitpunkt null. Auf dem statischen Plateau ist \(\mathbf w=-\mathbf u/c_\sigma\). Würde man dieses Verhältnis auch während der Schaltung vorgeben, verschwände die Summenspannung linear. Aus der Summe der beiden Bewegungsgleichungen folgte dann
\[
 (\rho_1-\rho_2/c_\sigma)\,\ddot{\mathbf u}=0.
\]
Mit \(\beta_\rho=\rho_2-c_\sigma\rho_1\) ist der Vorfaktor \(-\beta_\rho/c_\sigma\). Gerade bei \(\beta_\rho\ne0\) könnte dieser starre Ansatz also keinen Einschaltvorgang aus dem Ruhezustand tragen. Die zwei Operatoren in \eqref{eq:op5-elastic} lösen dieses Problem, indem die Summenspannung während der Schaltung die Änderung des Gesamtimpulses trägt \cite{op5-motion}.

Für die gewählte Verschiebung lauten die zugehörigen linearen mechanischen Beiträge
\begin{equation}
 \delta T^{00}=-\beta_\rho\nabla\cdot\mathcal L\mathbf b,\qquad
 \delta T^{0i}=\beta_\rho(\mathcal L\dot{\mathbf b})_i,\qquad
 \delta T^{ij}=-\beta_\rho\sigma_{ij}(\ddot{\mathbf b}) .
 \label{eq:op5-tensor}
\end{equation}
Zur Kontrolle kann man Dichte, Impuls und Spannung aus zwei verschiedenen Linearkombinationen der Verschiebungen lesen:
\begin{align}
 \rho_1\mathbf u+\rho_2\mathbf w
    &=\beta_\rho\mathcal L\mathbf b,
 &
 \mathbf u+c_\sigma\mathbf w&=\beta_\rho\ddot{\mathbf b}.
 \label{eq:op5-sum-identities}
\end{align}
Die Dichte ist minus die Divergenz der ersten Kombination, der Impuls ihre Zeitableitung. Die Spannungskomponenten des Quellentensors sind minus \(\sigma\) angewandt auf die zweite Kombination. Die Kontinuitätsgleichung kürzt so die Dichteableitung gegen die Impulsdivergenz. In der Impulsbilanz kürzt sich \(\beta_\rho\mathcal L\ddot{\mathbf b}\) gegen die Divergenz der Summenspannung. Für \(\mathbf b=\nabla S\) und \(K_L=\lambda+2\mu>0\), mit den Lamé-Koeffizienten \(\lambda,\mu\) der ersten Komponente, wird \(\mathcal L\mathbf b=K_L\nabla\Delta S\). Setzt man wie oben die Ruhedichte des Materials mit seiner Yukawa-Massendichte gleich, wird mit \eqref{eq:op5-density}
\begin{equation}
 \delta\mathcal J_h=-\frac{\delta T^{00}}{v}
                  =\frac{\beta_\rho K_L}{v}\Delta^2 S .
 \label{eq:op5-elas-source}
\end{equation}
Die zwei Laplace-Operatoren in \eqref{eq:op5-elas-source} kommen von der longitudinalen Elastodynamik und der Divergenz in der Dichte. Die Konstruktion mit einem direkten Higgsantrieb setzt \(\beta_\rho\ne0\) voraus. Bei \(\beta_\rho=0\) verschwinden in diesem Ansatz sowohl der mechanische Tensor als auch diese Quelle.

\subsubsection{Wo die Kopplung an das OP5-Feld eingeht}

Das reduzierte OP5-Modell setzt einen gerichteten massiven Vermittler mit Massenparameter \(m>0\) und ein skalares Vergleichsfeld \(h\) an. Für ein glattes Präparationsprofil \(H\) und verschwindende frühere Störungen gilt \(h=d_m H/m^2\) in Fourierdarstellung. Wir verwenden den Ansatz mit unverändertem freiem Einstein-Hauptteil. Für \(x=\omega^2\), \(q=|\mathbf k|^2\), \(a_\perp=|\mathbf k_\perp|^2\) und \(d=q-x\) hat der Antriebsoperator \(\mathcal P_{\rm dyn}\) das Symbol
\begin{align}
 d_m&=m^2+r_mq-x, &
 N_m&=\tfrac23(x^2-a_\perp x+a_\perp^2),\\
 \mathcal P_{\rm dyn}(\omega,\mathbf k)
 &=\frac{[Z(\Omega_0^2+c_h^2q-x)+\tfrac32g_0^2d]d_m-g_m^2N_m}{m^2}.
 \label{eq:op5-drive}
\end{align}
Die Koeffizienten und die gerichtete Quellenform sind vorgegeben. Eine in den Notizen geprüfte Wahl ist \(Z=1\), \(r_m=c_h^2=1/2\), \(g_0^2=g_m^2=3/100\), \(m^2,\Omega_0^2>0\). Sie dient hier als Rechenbeispiel. Liest man das Vergleichsfeld als Radialmodus des Higgsradius von W2, erzwingt W2 \(c_h=1\), bei minimaler Kopplung \(g_0=0\), und mangels gerichteten Vermittlers \(g_m=0\); \(c_h^2=1/2\) ist also eine Vergleichswahl, die W2 nicht trägt (Runde v4.10.0-r33 \cite{layers}). Die Kegelgrenze der Kausalität verlangt, dass beide Vermittler und der Higgsradius langsamer als das Licht laufen. Mit lichtschnellem Higgsradius ist die gerichtete Kopplung für keine Wahl der Vermittlergeschwindigkeiten verträglich (Runden v4.10.0-r47 und v4.10.0-r49 \cite{layers}). Im Vakuum begrenzen Laborgrenzen der Lorentz-Verletzung, die für Fermionen gemessen und auf den Higgs-Sektor übertragen sind, die Kopplung auf \(g_m^2\lesssim10^{-26}\). Der Beispielwert \(g_m^2=3/100\) ist dort ausgeschlossen und nur in einem Material denkbar, dessen Ruhesystem die Präparation ist (Runden v4.10.0-r50 und v4.10.0-r51 \cite{layers}; Anhang~\ref{sec:anhang-b-annahmen}). Auf ein Profil \(\delta H\) antwortet das reduzierte Modell bei passender Präparation auf den Antrieb \(\delta J_{\rm OP5}=\mathcal P_{\rm dyn}\delta H\) \cite{op5-channel}. Die Rechenschritte von hier bis zur Bilanz zweiter Ordnung sind in Anhang~\ref{sec:anhang-b} einzeln hergeleitet und mit Wolfram geprüft.

Die Wahl der zwei Quellenformen hat im reduzierten Modell eine konkrete Aufgabe. Mit \(p^\mu=(\omega,\mathbf k)\) ist \(d=p_\mu p^\mu\). Am Einstein-Kanal wird die identisch transversale Verbesserung \(S^{\mu\nu}=p^\mu p^\nu-\eta^{\mu\nu}d\) eingesetzt. In vier Dimensionen erfüllen ihre Spur und ihre doppelte Kontraktion
\[
 p_\mu S^{\mu\nu}=0,\qquad
 \operatorname{tr}S=-3d,\qquad S:S=3d^2 .
\]
Der spurgespiegelte Austauschterm der beiden Quellen ist proportional zu
\[
 \frac{S:S-(\operatorname{tr}S)^2/2}{d}=-\frac32 d .
\]
Nach Kürzung des Pols bleibt der lokale Beitrag, der in den Antriebsoperator \eqref{eq:op5-drive} eingeht. Auf derselben Verbesserung verschwindet auch die erhaltene Fierz--Pauli-Kontraktion \(S:S-(\operatorname{tr}S)^2/3\). Der gerichtete massive Kontrast in dieser Rechnung benutzt deshalb die eigens gewählte, zur Messachse transversale Quellenform mit dem Zähler \(N_m\). Die gemeinsame Herkunft beider Quellenformen aus einer physischen Wirkung bleibt zu zeigen \cite{op5-channel}. Ein erster Feldansatz gibt den Vermittlerterm mit eindeutigen Kopplungen exakt wieder, setzt aber einen zusätzlichen Skalar (Runde v4.10.0-r38 \cite{layers}). Für sich genommen ist der gerichtete Term ein gewöhnlicher Fierz--Pauli-Austausch einer gesetzten, gerichteten und erhaltenen Quelle; offen ist, ob Rahmen und Achse dieser Quelle zu einem Feld oder zur Auslese gehören (Runde v4.10.0-r45 \cite{layers}).

Die Materierechnung liefert \(\delta\mathcal J_h\). Der Vergleich verlangt \(\delta J_{\rm OP5}\). Wir können beide in deklarierter Quellen-Normierung gleichsetzen und ausrechnen, was daraus folgt. Eine Wirkung, die diese Gleichsetzung aus dem gemeinsamen Träger herleitet, liegt noch nicht vor. Am Higgs-Sektor von W2 folgt sie für den skalaren Kanal, wenn das Vergleichsfeld der Higgs-Radialmodus ist; dann fehlt aber der gerichtete Vermittler, an dem der Kontrast hängt (Runde v4.10.0-r33 \cite{layers}). Gleichung \eqref{eq:op5-match} markiert die zusätzliche Modellannahme; was sie trägt und was ohne sie bestehen bleibt, steht in Anhang~\ref{sec:anhang-b-annahmen}:
\begin{equation}
 \delta J_{\rm OP5}=\delta\mathcal J_h .
 \label{eq:op5-match}
\end{equation}
Mit \(D_0=\partial_t^2-\Delta\) wähle eine glatte, räumlich kompakte Funktion \(V\), die vor der Schaltung null ist, und setze
\begin{equation}
 \delta H=-D_0\Delta^2 V,\qquad
 S=-\frac{v}{\beta_\rho K_L}D_0\mathcal P_{\rm dyn}V .
 \label{eq:op5-profile}
\end{equation}
Die konstantkoeffizientigen Operatoren kommutieren. Beim Einsetzen in \eqref{eq:op5-elas-source} erhält man \(\delta\mathcal J_h=\mathcal P_{\rm dyn}\delta H\). Dieselbe Wahl hat eine zweite Konsequenz. Mit \(\mathbf b_*=(v/(\beta_\rho K_L))\nabla\mathcal P_{\rm dyn}V\) gilt \(\mathbf b=-D_0\mathbf b_*\). Folglich ist \(\delta T=-D_0 T_*\), wobei \(T_*\) der separat erhaltene Tensor aus \eqref{eq:op5-tensor} für \(\mathbf b_*\) ist.

Für die gewöhnliche linearisierte Gravitation verwenden wir harmonische Eichung. Ohne frühere Störung oder einlaufende freie Welle gehört zur Feldgleichung die folgende retardierte Lösung:
\begin{equation}
 D_0\bar\gamma_{\mu\nu}=16\pi G\,\delta T_{\mu\nu},
 \qquad \bar\gamma_{\mu\nu}=-16\pi G\,T_{*,\mu\nu} .
 \label{eq:op5-compact}
\end{equation}
Der gewöhnliche lineare Krümmungswechsel dieses präparierten mechanischen Beitrags verschwindet außerhalb seines räumlichen Trägers, auch während der Schaltung. Die Aussage benutzt den gemeinsamen Wellenfaktor der Quelle und die gewählten Anfangsdaten \cite{op5-motion}.

\subsubsection{Was auf derselben Messachse erhalten bleibt}

Ist damit auch der gerichtete Kontrast verschwunden? Auf dem statischen Plateau wird \(\delta H=\Delta^3 V\). Sei \(\xi=r_\perp^2/R^2\) die radiale Koordinate, \(\chi(\xi)\ge0\) ein glatter kompakter Zuschnitt, der nahe der Achse eins ist, und \(p(z)\) ein glattes Profil mit positivem Integral und Träger zwischen den Anschlüssen. Setze für \(\mathcal A\ne0\)
\begin{equation}
 V=-\frac{\mathcal A R^6}{230400}\,p(z)\chi(\xi)
       (\xi^5-\kappa\xi^6),\qquad
 \kappa=\frac{\int_0^\infty \xi^5\chi(\xi)\,d\xi}
              {\int_0^\infty \xi^6\chi(\xi)\,d\xi}.
 \label{eq:op5-axis-profile}
\end{equation}
Das Querintegral von \(V\) verschwindet je Längsschnitt.

Im flachen Kern kann man den Achsenterm Schritt für Schritt verfolgen. Mit der dimensionslosen Radialkoordinate \(\xi\) gilt \(\Delta_\perp \xi^j=4j^2\xi^{j-1}/R^2\). Dreimaliges Anwenden auf das Polynom aus \eqref{eq:op5-axis-profile} liefert \(230400(\xi^2-4\kappa\xi^3)/R^6\). Enthält ein Term von \(\Delta^3\) eine Längsableitung, bleiben ihm höchstens zwei radiale Laplace-Operatoren. Sein niedrigster radialer Grad ist daher \(\xi^3\). In der Nähe der Achse gelten für \(p(z)\ne0\) die normierten Formen
\begin{equation}
 \frac{\delta H}{-\mathcal A p(z)}=\xi^2+O(\xi^3),\qquad
 \frac{\Delta_\perp\delta H}{-16\mathcal A p(z)/R^2}=\xi+O(\xi^2).
 \label{eq:op5-normalized-axis}
\end{equation}
Die Profiländerung selbst ist auf der Achse null. Der massive Vergleich benutzt neben dem skalaren Feld \(h=d_mH/m^2\) das Feld \(f=\Delta_\perp H/m^2\), das aus demselben Profil mit dem transversalen Laplace-Operator gebildet wird. Nach \eqref{eq:op5-normalized-axis} ist auch sein Eintrag \(\delta f(0,z)\) auf der Achse null. Die zweite transversale Ableitung dieses Eintrags kann trotzdem von null verschieden sein \cite{op5-motion}.

\begin{figure}[!htbp]
\centering
\begin{tikzpicture}[>=Latex,x=43cm,y=13.4cm,font=\small]
  \draw[qs axis,->] (0,0)--(.215,0) node[right] {\(\xi=r_\perp^2/R^2\)};
  \draw[qs axis,->] (0,0)--(0,.211);
  \foreach \x in {.05,.1,.15,.2} {\draw[qs axis] (\x,-.004)--(\x,.004);}
  \node[anchor=north east] at (0,0) {\(0\)};
  \node[below=3pt] at (.1,0) {\(0{,}1\)};
  \node[below=3pt] at (.2,0) {\(0{,}2\)};
  \foreach \y in {.1,.2} {\draw[qs axis] (-.0012,\y)--(.0012,\y);}
  \node[left=4pt] at (0,.1) {\(0{,}1\)};
  \node[left=4pt] at (0,.2) {\(0{,}2\)};
  \draw[qs rust,very thick,dashed,domain=0:.19,samples=2] plot (\x,{\x});
  \draw[qs blue,very thick,domain=0:.19,samples=38] plot (\x,{\x*\x});
  \node[anchor=south east,qs rust] at (.177,.179) {\(\xi\)};
  \node[anchor=south west,qs blue] at (.163,.038) {\(\xi^2\)};
\end{tikzpicture}
\caption{Führende Terme nahe der Achse. Für \(p(z)\ne0\) beginnt das normierte Profil \(\delta H/(-\mathcal A p)\) mit \(\xi^2\), und sein normierter transversaler Laplace-Ausdruck \(\Delta_\perp\delta H/(-16\mathcal A p/R^2)\) mit \(\xi\). Damit verschwinden auf der Achse beide Werte, doch die zweite transversale Ableitung der Laplace-Antwort liefert den Gezeitenkontrast. Gezeigt werden die führenden Terme.}
\label{fig:op5-profil}
\end{figure}


In der geometrischen Optik ergibt sich für eine transversale Einheitsrichtung \(e\) auf derselben Achse der Gezeitenkontrast
\begin{equation}
 \Delta\mathcal I=-\frac12\int\partial_e^2(\delta f)(0,z)\,dz
                 =\frac{16\mathcal A}{m^2R^4}\int p(z)\,dz .
 \label{eq:op5-contrast}
\end{equation}
Der Ausdruck ist ungleich null, obwohl der in \eqref{eq:op5-compact} betrachtete gewöhnliche lineare Außenbeitrag kompakt ist. In der idealisierten Krümmungsauslese multipliziert eine gesetzte Differenz der effektiven Kopplungen der beiden Auslesekanäle diesen Kontrast. Erst mit einer solchen Kopplungsdifferenz und kalibrierten Anschlüssen lässt sich daraus eine Laufzeitdifferenz für einen Messaufbau bestimmen \cite{op5-axis,op5-motion}. Von der Ausleseseite gelesen ist die gerichtete Kopplung die kovariante Auslese des Gezeitenfelds quer zur Messachse eines Detektors mit Ruhesystem \(u\) und Achse \(n\) (Runde v4.10.0-r46 \cite{layers}). Rahmen und Achse wählt diese Auslese nicht selbst aus; sie sind eine Setzung des Detektors (Runde v4.10.0-r54 \cite{layers}).

Der Kontrast \(\Delta\mathcal I\) ist ein Integral einer Krümmungskomponente des Vergleichsfelds für eine feste Polarisation. Ein statischer konformer Term geht auf der Achse als integrierte Randableitung ein, und ein rein longitudinaler Anteil hat keine lineare Krümmung. Die Quelle kann so an den Enden identische lokale Vergleichsfelder liefern und auf dem Weg trotzdem den gerichteten Kontrast ändern \cite{op5-axis}. Für die dynamisch präparierte Profilfamilie liefert die Azimutalmittelung über ein rotationssymmetrisch gewichtetes, achsenparalleles Bündel mit konstantem Querschnitt im flachen Kern den Faktor
\begin{equation}
 1-36\kappa\langle\xi\rangle
 \label{eq:op5-beam}
\end{equation}
zum auf der Achse berechneten reinen Riemann-Kontrast. In dieser Mittelung verschwinden die Längsableitungsbeiträge nach der Integration in \(z\). Der Faktor ist eine Kontrolle der Näherung am statischen Plateau. Bei Beugung, wechselndem Querschnitt oder Durchlaufen der Strecke während des Umschaltens ist die optische Übertragung eine weitere Rechnung \cite{op5-motion,op5-axis}.

\paragraph{Planare Kontrolle.}
Beim planaren statischen Vergleich mit dem Faltungskern \(K_m(z)=m e^{-m|z|}/2\) und einer im Inneren der Strecke getragenen Dichtedifferenz \(\delta\rho\) gilt für \(f=K_m*\delta\rho\)
\begin{equation}
 \int_0^L f(z)\,dz
   =\int_0^L\delta\rho(z)\,dz-\frac{f(0)+f(L)}{m}.
 \label{eq:op5-planar-check}
\end{equation}
Die Faltung erhält das Integral über die ganze Gerade. Außerhalb des Trägers sind ihre beiden Ausläufer exponentiell und ihre Integrale werden durch die beiden Endwerte bestimmt. Bei gleicher integrierter Dichte und gleichen Endwerten in den beiden Stellungen bleibt damit aus diesem planaren Beitrag keine führende Differenz. In der axialen Konstruktion von \eqref{eq:op5-contrast} misst dieselbe geometrische Strecke eine transversale Krümmungsableitung. Gleiches Querintegral und gleiche lokale Endwerte schließen diesen Kontrast nicht aus. Das ist die Kontrolle, die den Wechsel des Profils mit festem Messweg motiviert \cite{op5-axis}.

\subsubsection{Warum die zweite Ordnung eine eigene Bilanz braucht}

\paragraph{Arbeit beim Schalten.}
Die innere Feder überträgt Impuls zwischen den beiden Materialien. Ihre vorgeschriebene Ruhelage \(\mathbf d\) ist zugleich eine Arbeitsvariable. Setze \(\mathbf R_{\rm F}=\mathbf u-\mathbf w-\mathbf d\), \(\sigma_1=\sigma(\mathbf u)\) und \(\sigma_2=c_\sigma\,\sigma(\mathbf w)\). Für die zur zweiten Ordnung beitragende mechanische Energiedichte gilt
\begin{align}
 \partial_t e_{\rm mech}&=\nabla\cdot\mathbf F+P_{\rm Schalt},\\
 \mathbf F&=\sigma_1\dot{\mathbf u}+\sigma_2\dot{\mathbf w},
 &P_{\rm Schalt}&=-k\,\mathbf R_{\rm F}\cdot\dot{\mathbf d}.
 \label{eq:op5-drive-energy}
\end{align}
Die Vorzeichen folgen aus der Federenergie \(k|\mathbf R_{\rm F}|^2/2\). Beim Differenzieren hebt sich die Arbeit der beiden inneren Federkräfte gegen den Anteil mit \(\dot{\mathbf u}-\dot{\mathbf w}\) auf. Übrig bleibt der von außen vorgegebene Anteil mit \(\dot{\mathbf d}\). Ein ortsfester Speicher mit \(\partial_t e_{\rm res}=-P_{\rm Schalt}\) gleicht die Arbeit im Volumen aus. Aus der Summe folgt auf dem Plateau noch der Dichtebeitrag \(\nabla\cdot\int\mathbf F\,dt\). Sein Volumenintegral verschwindet für kompakten Fluss. Höhere Momente können erhalten bleiben \cite{op5-feder}.

\paragraph{Welche Verschiebung am Plateau fehlt.}

Die lineare Gegenbewegung ist ein Befund für die präparierte Quelle. Um die zweite Ordnung zu prüfen, betrachten die Forschungsnotizen dieselben beiden Hauptkomponenten auf dem Plateau mit einem Saint-Venant--Kirchhoff-Gesetz und ausschließlich innerer Paarfeder. Sei \(\mathsf H=\nabla_X\mathbf u_1\) die symmetrische erste Deformation im Referenzchart, mit \(\mathbf u_1=\nabla_X\varphi\) und \(\mathbf w_1=-\mathbf u_1/c_\sigma\). Für die zweite Summenverschiebung \(S_2=\mathbf u_2+c_\sigma\mathbf w_2\) hat die Summe der Piola-Tensoren den Koeffizienten
\begin{align}
 (P_1+P_2)^{(2)}&=\mathsf C\,\varepsilon(S_2)
                     +(1+1/c_\sigma)N(\mathsf H),\\
 N(\mathsf H)&=\tfrac{\lambda}{2}\operatorname{tr}(\mathsf H^2)I
                 +3\mu\mathsf H^2+\lambda(\operatorname{tr}\mathsf H)\mathsf H .
\end{align}
Für \(\mu>0\), \(3\lambda+2\mu>0\) und eine nichttriviale kompakte longitudinale erste Verschiebung folgt
\begin{equation}
 \int\operatorname{tr}N(\mathsf H)\,d^3X
 =\left(\tfrac52\lambda+3\mu\right)\int|\mathsf H|^2\,d^3X>0 .
 \label{eq:op5-second}
\end{equation}

Bei kompaktem \(S_2\) erzwingt die statische Summengleichung nach Testen mit \(X\) eine verschwindende integrierte Spur. Der lineare Dehnungsterm integriert sich dabei zu null. Der positive quadratische Beitrag widerspricht dieser notwendigen Bedingung. Damit scheitert eine kompakte zweite Summenverschiebung unter den angegebenen Material- und Halterungsannahmen \cite{op5-elastic-limit}.

Die Spuridentität lässt sich im selben Referenzchart prüfen. Für die Hesse-Matrix \(\mathsf H\) von \(\varphi\) mit kompakter erster Verschiebung ist
\[
 (\Delta\varphi)^2-\operatorname{tr}(\mathsf H^2)
 =\nabla_X\cdot\big[(\Delta\varphi)\nabla_X\varphi
                    -\mathsf H\nabla_X\varphi\big].
\]
Der Randterm verschwindet. In der integrierten Spur von \(N\) lassen sich daher \((\operatorname{tr}\mathsf H)^2\) und \(\operatorname{tr}\mathsf H^2\) zusammenfassen. Damit erhält man den Koeffizienten aus \eqref{eq:op5-second}. Das Argument prüft eine mechanische Summenverschiebung und verwendet keine optische Auslese. Es macht auch deutlich, welche Größen getrennt bilanziert werden müssen. Der zusätzliche Term im Restmassentransport ist
\begin{equation}
 -\nabla\cdot(\rho_1\mathbf u_2+\rho_2\mathbf w_2).
 \label{eq:op5-second-density}
\end{equation}
Die statische Summengleichung in \(X\) bestimmt die Kombination mit \(c_\sigma\). Der gravitative Beitrag aus \eqref{eq:op5-second-density} braucht die mit den Ruhedichten gewichtete Kombination. Für den räumlichen Stress ist außerdem der Cauchy-Tensor am selben Raumpunkt zu bilden. Die Federkräfte setzen dort jeweils an unterschiedlichen Endpunkten an und benötigen eine Bilanz des Verbindungsstresses. Bei nichtzentralen Kräften kommt die Drehmomentbilanz der Halterung hinzu \cite{op5-elastic-limit}.

Der vorhandene Reservoirabgleich kompensiert die bereits berechneten direkten Dichteanteile auf dem Plateau. Der vollständige gewöhnliche Beitrag zweiter Ordnung bleibt zu bestimmen \cite{op5-reservoir,op5-elastic-limit}.

\subsubsection{Wie sich eine Laufzeit von der Anschlussphase trennen ließe}
\label{sec:op5-calibration}

Die Krümmung in \eqref{eq:op5-contrast} ist eine Größe im Modell. Wie sähe eine Auslese aus, die sie von einer bloßen Änderung der Anschlüsse trennen könnte? Betrachte dazu eine kohärente, verlustfreie Zweimodenstrecke in der linearen Einweg-Näherung. Nach dem Ausklammern der gemeinsamen Grundlaufzeit seien \(M_A=LF\) und \(M_B=LUF\). \(L\) und \(F\) sind invertierbare Empfangs- und Sendematrizen. \(U\) ist der unitäre Streckenkontrast zwischen den Profilstellungen. Wenn die beiden Enden in beiden Stellungen dieselben Matrizen haben, gilt
\begin{equation}
 M_A^{-1}M_B=F^{-1}UF .
 \label{eq:op5-ab-calibration}
\end{equation}
Die Eigenphasen sind dann diejenigen der Strecke. Diese Aussage benutzt die Ähnlichkeit und braucht keine diagonalen Anschlüsse. Falls sich die Enden mit dem Profil ändern, kann ein profilsynchroner Phasenterm am Sender denselben Frequenzgang wie ein mit der Laufzeitmatrix kommutierender Zusatz auf der Strecke erzeugen. Der Quotient allein entscheidet in diesem Fall nicht über den Ort des Effekts \cite{op5-budget}.

Ein unabhängiger Nullweg mit bekanntem Einheitskontrast würde im Idealfall die gleichen aktuellen Endmatrizen teilen. Für relative Abweichungen \(S,R\) am Empfangs- und Sendeende lässt sich diese Forderung direkt schreiben. Ist \(C=LF\) die Nullmessung und \(M=LSURF\) die Messung, so ist \(C^{-1}M\) für invertierbares \(S\) ähnlich zu \(U(RS)\). Reziproke Abweichungen mit \(RS=I\) erhalten das Streckenspektrum. Bei \(RS=c_0 I\), \(c_0\ne0\), erhalten beide Eigenwerte denselben Faktor, sodass ihre korrekt entfaltete Phasendifferenz gleich bleibt. Alternativ könnte eine im Messzweig unabhängig begründete Nullstellung \(M_0=LSRF\) die gleichen relativen Anschlüsse wie die Messung verwenden. Dann gilt \(M_0^{-1}M=F^{-1}R^{-1}URF\). Die Voraussetzung ist in beiden Fällen eine kohärent geteilte Endpunktantwort. Die axiale Nullstellung des Feldvergleichs liefert dafür ein Kriterium an das Profil. Sie bestimmt noch nicht den Abgleich eines physischen Messanschlusses \cite{op5-nullpfad,op5-axis}.

Auch bei abgeglichenen Enden bleibt eine Anforderung an die Stabilität der Phase. Schreibe die fehlerbehafteten Matrizen in der Feldbasis als \(\widehat M_A=L(I+N_A)F\) und \(\widehat M_B=L(U+N_B+D_{\rm E})F\), mit \(D_{\rm E}\) als Änderung der Endpunktantwort in Stellung B. Für invertierbares \(I+N_A\) ist der Messquotient \(\widehat M_A^{-1}\widehat M_B\) ähnlich zu \(Q=(I+N_A)^{-1}(U+N_B+D_{\rm E})\). Ihr Fehler gegenüber \(U\) ist exakt
\begin{equation}
 Q-U=(I+N_A)^{-1}(N_B-N_AU+D_{\rm E}) .
 \label{eq:op5-phase-error}
\end{equation}
Bezeichne \(r\) eine Schranke für \(\|Q-U\|_2\). Für normales \(U\) liegt jeder gestörte Eigenwert höchstens im Abstand \(r\) zu einem idealen Eigenwert. Sind die Kreise disjunkt, bleibt die Zuordnung erhalten. Für \(r<1\) beträgt die Winkelabweichung jedes zugeordneten Eigenwerts höchstens \(\arcsin r\). Über zwei Frequenzen mit Abstand \(\Delta\omega\) erhält man für die Sekante der differentiellen Phase nach korrekter Entfaltung die Fehlergrenze \(4\arcsin r/\Delta\omega\). Die vier Beiträge stammen von zwei Eigenphasen an jeweils zwei Frequenzen \cite{op5-budget}. Die Kalibrieraufgabe erhält damit eine präzise Form. Das am Modell vorhergesagte Signal muss sich sowohl gegen die verbleibende Endpunktänderung als auch gegen die Phasenfehler behaupten.

Die Fallstudie zeigt damit, welche Schritte sich gemeinsam rechnen lassen. Die W2-Massenparameter bestimmen die radiale Materiequelle. Für das präparierte Profil sind gewöhnliches lineares Außenfeld und gerichteter Achsenkontrast getrennt berechenbar. Der gewöhnliche Teil folgt allein aus der linearisierten Einstein-Gleichung. Der gerichtete Teil hängt an der Kopplungsannahme \eqref{eq:op5-match} und an einem gerichteten, unterluminalen massiven Vermittler, den W2 nicht enthält. Mit dem lichtschnellen Higgsradius von W2 verletzt die gerichtete Kopplung die Kegelgrenze, und im Vakuum bleibt nur \(g_m^2\lesssim10^{-26}\); ein beobachtbarer Kontrast verlangt ein Material. Die native Massenauslese, eine Wirkung für Vermittler und Kopplung und die kalibrierte Auslese bleiben weitere Schritte \cite{op5}; OP5 bleibt offen. Den Stand der Vermittlersuche fasst Anhang~\ref{sec:anhang-b-annahmen} zusammen, die Einordnung unter die übrigen offenen Probleme Runde v4.10.0-r55 \cite{crossmap}.

\section{Woran die Idee sich prüfen lässt}
\label{sec:tests}

Die Aufgabe eines gemeinsamen Trägers ist erst dann erfüllt, wenn dieselben Zustände, Symmetrien und Kopplungen in den aufeinanderfolgenden Rechnungen wieder auftauchen. Die Beispiele machen diese Anforderung an drei Stellen konkret. Sie zeigen auch, wie sich ein negativer Befund zur jeweils benutzten Annahme zurückverfolgen lässt.

\subsection{Was die Feldauslese erhält}

Beim HC-Feldvergleich sind die Eingaben eine Feldkonfiguration und die deklarierte Involution \(J_{\mathrm B}\). Abbildung \ref{fig:hc-feld} zeigt den Test am lichtartigen Eingang. Das berechnete \(s'\) sagt, welche Feldklasse nach der Projektion erreicht wird. Ein Modell, das aus demselben Projektor wieder ein lichtartiges Feld gewinnen soll, muss angeben, durch welche anschließende Operation die beiden Bedingungen \(s'=0\) und \(\mathbf z'\ne0\) erfüllt werden. Die Feldinvariante macht diese Forderung unabhängig von der Wahl eines Lorentz-Beobachters prüfbar. Die Identifikation eines quantisierten Lichtzustands setzt weitere Struktur voraus \cite{feldkern,feldsektoren}.

Auch die Dynamik der Auslese hat einen eigenen Test. Für ein vorgegebenes Paar \((T(t),\allowbreak G(t))\) genügt ein \(u_{\ker}\in\ker T\) mit \((\dot T+TG)u_{\ker}\ne0\), um eine autonome lineare Entwicklung der ausgelesenen Variablen auszuschließen. Die Zustände \(u\) und \(u+u_{\ker}\) ergeben dieselbe Auslese und verschiedene Ableitungen. Die Gleichheit der momentanen Beobachtungsdaten entscheidet so noch nicht, ob zwei Zustände dieselbe Antwort auf eine Störung haben. Diese Prüfung lässt sich an einer konkreten Realisierung durchführen, sobald ihre Auslese und ihr Generator feststehen \cite{physics}.

\subsection{Was bei den Massen getrennt festgelegt werden muss}

Im Gram-Vergleich sind die Beträge der paarweisen Überlappungen und die Phase getrennte Eingaben. Man kann die Phase ändern und alle Beträge festhalten. Die Normrelation bleibt erhalten, während die Determinante und das Spektrum sich ändern. Abbildung \ref{fig:koide-phase} zeigt den Rangverlust an einer konkreten Stelle dieser Familie. Der Versuch, die drei Leptonenmassen aus der Überlappung allein zu bestimmen, lässt daher einen ausrechenbaren Teil des Spektrums frei.

Welche Prüfung folgt daraus? Eine vorgeschlagene Zustandspräparation muss zuerst unabhängig die zugehörige Gram-Matrix liefern. Falls sie Rang zwei hat, liefert diese Wurzelregel eine verschwindende Masse. Den Randfall kann man angeben, bevor man eine Skala wählt. Andererseits liefert ein positiv definites Beispiel zunächst drei positive Wurzeln. Für einen Vergleich mit physikalischen Massen braucht man danach die Identifikation des nativen Operators, die Auswahl der Phase und eine unabhängig bestimmte Normierung. Ein erfolgreicher Test der Normrelation beantwortet in diesem Aufbau genau eine der drei Fragen \cite{physics,companion}.

\subsection{Was ein OP5-Vergleich messen müsste}

Bei OP5 ist die Eingabe das präparierte Profil aus \eqref{eq:op5-profile}. Abbildung \ref{fig:op5-modell} zeigt, wie daran zwei Antwortwege anschließen. Man kann die Erhaltung des mechanischen Tensors prüfen und für den gerichteten Weg den Gezeitenkontrast auf dem Plateau berechnen. Der gewöhnliche Außenbeitrag \eqref{eq:op5-compact} und der Achsenkontrast \eqref{eq:op5-contrast} sind damit am selben Profil getrennt zugänglich. Das Nahachsenverhalten in Abbildung \ref{fig:op5-profil} zeigt außerdem, warum ein auf der Achse verschwindender Profilwert den Gezeitentest offenlässt. Die Auslese benutzt Ableitungen des Profils.

Eine spätere Messung braucht für beide Antwortwege einen Bezug zur realen Präparation. Die Beiträge des Materials, der realen Halterung und der Ansteuerungsenergie bestimmen den gewöhnlichen Hintergrund. Die Empfindlichkeit der Anschlüsse und die effektive Kopplungsdifferenz bestimmen, ob sich der gerichtete Kontrast in einer Übertragungsdifferenz wiederfindet. Die integrierte Spur in \eqref{eq:op5-second} lokalisiert ein konkretes Problem der linearen Kompensation in der nächsten Materialordnung. Für diese Ordnung ist ein vollständiger Tensor am räumlichen Ort der Messung zu bilanzieren \cite{op5-elastic-limit,op5-budget}. Der angegebene Grenzfall liefert damit eine überprüfbare Anforderung an eine Apparatur.

\subsection{Was an einer einzigen Zahl hängt}
\label{sec:one-number}

Die Zahl stammt aus der Zählung der Zustände eines Horizonts in der
Schleifenquantengravitation. Eine Horizont-Zählregel ist eine Gleichung
\[
 \sum_j g(j)\,e^{-2\pi\gamma\sqrt{j(j+1)}}=1,\qquad j=\tfrac12,1,\tfrac32,\dots,
\]
deren Lösung \(\gamma\) den Barbero-Immirzi-Parameter festlegt; \(g(j)\) ist das
Gewicht des Spins \(j\). Der Materie-Frame M zählt jeden Spin mit dem Gewicht
\(g(j)=2\), nach der Regel von Domagala und Lewandowski. Das ergibt \(\gamma_M=0{,}237533\). Der
Tachyon-Frame T zählt nur \(j=\tfrac12\), ebenfalls mit dem Gewicht 2, nach der
Regel von Ashtekar, Baez, Corichi und Krasnov. Das ergibt
\(\gamma_T=\ln2/(\pi\sqrt3)=0{,}127384\) (Runden v4.10.0-r52 und v4.10.0-r53
\cite{layers}). Der Quotient ist der Frame-Faktor
\(\gamma_L=\gamma_M/\gamma_T=1{,}8647\). Schreibt man \(\gamma_L=\cosh2\eta\)
und \(\beta=\tanh\eta\), gehört dazu die Geschwindigkeit \(\beta=0{,}5494\)
(Rapidität \(\eta=0{,}6175\)). OP1, das erste offene Problem des Companion,
ist die Frage, warum \(\beta\) diesen Wert hat; gleichwertig: warum M und T
nach verschiedenen Regeln zählen. In der Beobachterlesart ist M ein
Beobachter, der sich relativ zum Ruhesystem eines Feldes mit parallelem
elektrischem und magnetischem Teil bewegt. OP1 fragt dann, warum er sich mit
\(\beta=0{,}5494\) bewegt; \(\gamma_L\) ist in dieser Lesart das Verhältnis
der Energiedichten des Feldes, gemessen von M und im Ruhesystem (Runde
v4.10.0-r29 \cite{companion}).

Seit dem 26. November 2026 ist die Zählregel als \emph{Postulat} angenommen (Runde v4.10.0-r67 \cite{countingpostulate}), im selben Sinn
wie die Gleichberechtigung der Frames: benannt, nicht hergeleitet. Es hat drei Teile. Im Materie-Frame zählt jeder Spin mit dem Gewicht 2
(die beiden äußersten Zustände \(m=\pm j\)); die Fläche eines Durchstoßpunkts ist proportional zu \(\sqrt{j(j+1)}\); und der
Tachyon-Frame zählt nur \(j=\tfrac12\), wobei ein solcher Punkt ein Bit trägt. Mit \(k=2j\) folgt daraus eine einzige Gleichung für den
Frame-Faktor,
\[
 \sum_{k\ge1}2^{\,1-\gamma_L\sqrt{k(k+2)/3}}=1,\qquad \gamma_L=1{,}8646997646,\qquad \beta=0{,}5494056410 .
\]
Damit gibt es eine Rechnung, die zeigt, woher die Zahl kommt; es fehlt der Grund, warum die Zählung diese Form hat. Ersetzt man die Fläche
durch \(j+\tfrac12\), gibt dieselbe Regel exakt \(\beta=1/\sqrt3\). Der Wert ist nicht an Daten angepasst, die Regel wurde aber in
Kenntnis der Datenfenster gewählt. OP1 bleibt offen und lautet jetzt: Leite das Zählpostulat her.

Dieselbe Runde und Runde v4.10.0-r66 halten fest, was das Bild zweier Auslesen dazu beiträgt. Liest man die Hopf-Sphäre als Himmelskugel
von Lichtrichtungen (eine gesetzte Lesart), dann ist \(\beta\) der Überlapp des Hopf-Zustands mit dem Gegenpol des Photonpols, und der
Boost, der beide gegenüberstellt, hat die Geschwindigkeit \(\beta\); Winkel-Lesart und Geschwindigkeits-Lesart sind dann dieselbe
Aussage. Das gilt für jeden Wert von \(\beta\) und leitet keinen her. Die naheliegende Schließungsbedingung dieses Bildes, dass die
Übergangswahrscheinlichkeiten zwischen den zwei Auslesen über alle Spins zusammen eins ergeben, lautet \(\sum_j2\beta^{4j}=1\) und
erzwingt \(\beta=1/\sqrt3\), nicht den gemessenen Wert (Test H165, widerlegt, überwiegend nicht blind). Eine Sichtung zählt 37 frühere
Versuche, die Zahl herzuleiten; keiner gelingt.

Das Zeichen \(\gamma\) hat in diesem Abschnitt drei Bedeutungen. Erstens ist
\(\gamma_L\) ein Frame-Faktor; er hat die Form eines Lorentz-Faktors,
\(\cosh2\eta\), und ist größer als eins. Zweitens sind \(\gamma_M\),
\(\gamma_T\) und das \(\gamma\) des Holst-Terms unten Werte des
Barbero-Immirzi-Parameters. Drittens ist derselbe Wert in geeigneten Einheiten
die Dimension einer selbstähnlichen Menge; \(\gamma_1\) am Ende des Abschnitts
ist eine solche Dimension.

Die bedingten Zahlenwerte des Companion gehen auf diese eine
dimensionslose Zahl zurück. In die
offenen Probleme gelangt sie über fünf getrennte Gleichsetzungen: einen
Spektralversatz, die Restbiegung der Leptonphase, \(|U_{e2}|^2=\beta^2\) in
der Neutrinomischung, \(\Omega_M=\beta^2\) für den Materieanteil und
\(\cot\theta_W=\gamma_L\). Vier davon wurden erst gewählt, als die Daten schon
bekannt waren (Bemerkung v4.10.0-r30.6 \cite{companion}). Jedes Problem behält daneben einen eigenen Rest, der von der
Zahl nicht abhängt \cite{rootmap}.

Wie viel die Treffer wiegen, prüft eine Zufallsabschätzung nach einem vorab
festgelegten Verfahren. Der Katalog umfasst sechzehn Ziele, acht Formen und
ein Fenster von \(2\sigma\). Das Paar aus der Regel von M und der Regel von T
trifft dreimal, als einziges der sechsunddreißig möglichen Paare von
Zählregeln. Für ein zufälliges \(\gamma_L\) liegt die Wahrscheinlichkeit für
mindestens drei Treffer bei \(0{,}018\) bis \(0{,}023\). Über alle
sechsunddreißig Paare erwartet man daher \(0{,}6\) bis \(0{,}8\) Paare mit
mindestens drei Treffern. Nur wenn das Paar ohne Blick auf die Daten
festgelegt worden wäre, läge die Wahrscheinlichkeit bei etwa \(2\,\%\).
Solange das Paar gewählt und nicht hergeleitet ist, sind die
Treffer mit Zufall verträglich. Nachträglich angemerkt: Die beiden getroffenen
Ziele nahe \(0{,}30\) sind stark korreliert; zwei der drei Treffer sind daher
fast einer (Bemerkung v4.10.0-r30.5 \cite{chance,companion}).

Die Zahl selbst beruht darauf, zwei Frames verschiedene Zählregeln
zuzuordnen. Das ist gleichwertig mit einem frame-abhängigen Flächenquant.
Ein unitär wirkender Boost ändert das Spektrum des Flächenoperators aber
nicht, nur die Verteilung seiner Eigenwerte \cite{rovelli}. Ein
frame-abhängiges Quant bräuchte daher eine nicht-unitäre Frame-Abbildung.
Ein Vermittler, der alle fünf Gleichsetzungen aus einer Kopplung erzeugt,
ist im Companion nicht gefunden. Die Physik-Ausgabe von v5.3.0 behandelt dies in Abschnitt~7.7 \cite{physics}.
Auch ein Holst-Term in der Gravitationswirkung liefert die Abbildung nicht:
Lorentz-Transformationen, selbst ein komplexer Boost, lassen den
Barbero-Immirzi-Parameter fest, und nur eine Dualitätsdrehung der Krümmung um
einen Winkel \(\theta\) ändert ihn nach \(\arctan\gamma\mapsto\arctan\gamma+\theta\).
Kein im Companion ausgezeichneter Winkel bildet den Wert des Materie-Frames
auf den des zweiten Frames ab (Runde v4.10.0-r31 \cite{holst,layers}). Mehr noch: Jede Dualitätsdrehung,
reell oder komplex und in beliebiger Folge, erhält die Quadratsumme der
Einstein- und Holst-Koeffizienten. Für reelle Ergebnisse unterscheiden sich die
Flächenquanten beider Frames dann um
\(\sqrt{(1+\gamma_T^2)/(1+\gamma_M^2)}\approx0{,}98\), nicht um den bei gleicher
Gravitationskonstante nötigen Faktor \(\gamma_T/\gamma_M\).
Auch die nicht-unitäre Frame-Abbildung \(J_{\mathrm B}\) liefert ihn nicht (Runde v4.10.0-r52 \cite{layers}): Sie ist ein Automorphismus der Lorentz-Algebra, und eine
Konjugation, unitär oder nicht, lässt jedes Spektrum fest. Liest man die Fläche über die transformierte Simplizitätsbedingung ab, wird \(\gamma\)
quer zur Frame-Achse zu \(-1/\gamma\) und bleibt entlang der Achse; das ergibt weder \(\gamma_T\) noch \(\gamma_L\).
In den Einheiten \(r_j=e^{-2\pi\sqrt{j(j+1)}}\) ist jede Zählregel die Ähnlichkeitsgleichung einer selbstähnlichen Familie auf der Geraden mit
der Dimension \(\gamma<1\) (Runde v4.10.0-r53 \cite{layers}). Eine neuere, noch unbegutachtete Dimensionsformel über Entropieraten, die exakte Überlappungen zulässt \cite{openai148}, zeigt, was das
Gewicht 2 der Materie-Regel verlangt: Die beiden Extremzustände \(|j,\pm j\rangle\) müssen als verschiedene Abbildungen wirken. Fallen sie zusammen
oder wird nur einer zugelassen, sinkt der Wert auf \(\gamma_1=0{,}1616\), und die Regel des zweiten Frames fällt auf die Dimension null. Das Gewicht 2
verlangt also eine Auslese, die die beiden Orientierungen einer Punktierung unterscheidet, ohne eine auszuwählen; warum sie das tut, ist nicht
hergeleitet.

\subsection{Welche Messungen die Zahl prüfen}
\label{sec:data}

Aus der Zahl \(\beta\) und ihren Gleichsetzungen folgen Werte, die Messungen prüfen können. Der Companion hält für jede eine Vorab-Notiz
mit Prüfgröße, Prüfdatum und Entscheidungskriterium fest, bevor die entscheidenden Daten vorliegen \cite{predictions}. Die Tabelle
nennt die Werte und ihren Stand am 23.~November 2026. Ein Wert gilt erst als geprüft, wenn eine Messung nach dem festgelegten Stichtag
vorliegt; frühere Daten werden nur zum Vergleich angegeben.

\begin{center}
\small
\begin{tabular}{@{}p{.07\linewidth}p{.25\linewidth}p{.29\linewidth}p{.29\linewidth}@{}}
\hline
\textbf{Name} & \textbf{Vorhersage} & \textbf{entscheidende Messung} & \textbf{Vergleich heute}\\
\hline
H9 & \(\sin^2\theta_{12}=\beta^2/\cos^2\theta_{13}=0{,}3086\) & JUNO mit Fehler \(\le0{,}004\) & JUNO 2026: \(0{,}3036\pm0{,}0064\), Abstand \(+0{,}8\sigma\) \cite{juno2026}\\
H96 & \(\Omega_M=\beta^2=0{,}30185\) & erstes \(\Omega_m\) aus DESI DR3 (nur BAO) oder Euclid; widerlegt bei mehr als \(2\sigma\) & DESI DR2: \(0{,}2975\pm0{,}0086\), Abstand \(+0{,}5\sigma\) \cite{desi-dr2}\\
H102 & \(m_\tau=1777{,}134\)\,MeV aus der Restbiegung der Leptonphase & erste PDG-Ausgabe nach dem 13.~November 2026; widerlegt bei mehr als \(2\sigma\) & PDG 2026: \(1776{,}93\pm0{,}09\)\,MeV, Abstand \(+2{,}3\sigma\); zählt nicht, weil vor der Vorhersage \cite{pdg}\\
H112 & \(\sin^2\theta_{12}\cos^2\theta_{13}=\Omega_m\), ohne den Wert von \(\beta\) & die Messungen von H9 und H96 zusammen; widerlegt bei mindestens \(3\sigma\) & JUNO gegen DESI DR2: \(-0{,}1\sigma\)\\
H15 & \(m_W=80{,}3615\pm0{,}0018\)\,GeV & erste Weltmittelung nach dem 21.~Oktober 2026 & PDG 2026: \(80{,}3625\pm0{,}0077\)\,GeV; zählt nicht, weil vor der Vorhersage\\
\hline
\end{tabular}
\end{center}

Am 23.~November 2026 ist keiner dieser Tests entschieden. Drei Einschränkungen gehören zu jedem Treffer. Erstens hängen alle Werte bis
auf H112 an der Zuordnung der Zählregeln, also an OP1; die Zufallsabschätzung in Abschnitt~\ref{sec:one-number} gilt deshalb auch hier.
Zweitens sind H9 und H102 nach dem Sehen der älteren Daten gefunden worden; blind sind sie nur gegenüber künftigen Messungen. Drittens
prüft H112 eine schwächere Aussage: dass an beiden Stellen dieselbe Zahl steht, gleich welche. Er widerlegt erst ab einem Unterschied
von etwa \(0{,}022\) in \(\beta^2\). Weitere Vorhersagen zum Mischungswinkel \(\theta_{23}\) warten auf die nächste globale Auswertung
der Neutrinodaten \cite{predictions}.

\subsection{Welche Rechnung eine Annahme mitführt}

Für eine Nachrechnung wählt man eine Gleichung und schreibt ihre Eingaben auf. Beim kompakten gewöhnlichen Außenbeitrag sind das der lineare erhaltene mechanische Tensor, der Wellenfaktor des Profils und die retardierten Daten. Die gerichtete Auslese benutzt dieselbe Präparation, dazu den gewählten OP5-Antriebsoperator und die Kopplungsannahme in \eqref{eq:op5-match}. Die Gleichsetzung der Quellen ist der Übergang, durch den beide Wege Antworten auf dieselbe Materieänderung werden \cite{op5-motion,op5}.

QAM hält zu einer Formel ihre Darstellung, ihre Voraussetzungen und die noch offenen Übergänge für ihre Anwendung fest. Die im Companion protokollierten symbolischen Kontrollen gehören jeweils zu den durchgeführten Rechnungen. Eine Änderung der physikalischen Kopplung würde die Anwendung auf das gemeinsame Modell berühren. Die Erhaltungsrechnung für den mechanischen Tensor hat ihre eigenen, im Fallbeispiel sichtbaren Voraussetzungen \cite{qam,op5-motion}.

\paragraph{Wolfram-Rechenkontrolle.}
Für Fassung 0.1.0 liefen in Wolfram Language 15.0.1 insgesamt 80 exakte Prüfgruppen zum HC, zur Feldalgebra, zu Koide und zum geometrischen Vergleich. Die separate Dixon-Rechnung reproduzierte die Kommutantendimensionen sechs und null für den komplex-linearen und den antilinearen Anteil. Anhang~\ref{app:wolfram-output} stellt die Wolfram-Eingaben und ihre Originalausgaben einander gegenüber. Das Wolfram-Begleitarchiv enthält die Eingabeskripte, die protokollierten Ausgaben, den Prüfbericht und den geprüften Stand. Die Rechnungen kontrollieren die angegebenen mathematischen Identitäten und Testfälle unter ihren jeweiligen Voraussetzungen. Sie betreffen nur den Stand der Fassung 0.1.0. Teile, die danach hinzugekommen oder erweitert worden sind, decken sie nicht ab, darunter die Abschnitte~\ref{sec:octant-sign} und \ref{sec:one-number}, die Erweiterungen in Abschnitt~\ref{sec:open-points} und Anhang~\ref{sec:anhang-b}. Für die Herleitungen in Anhang~\ref{sec:anhang-b} gibt es eine eigene, enger gefasste Prüfung (Abschnitt~\ref{sec:anhang-b-pruefung}). Die neuen Zahlen in Abschnitt~\ref{sec:one-number} und das Ergebnis zur Auslese in Abschnitt~\ref{sec:open-points} belegt Anhang~\ref{app:new-calcs} mit den Originalausgaben; für die übrigen Aussagen gelten die Belege der jeweils genannten Companion-Runden \cite{companion,layers}.

Für die Vereinheitlichungsfrage ergibt sich so eine Folge von konkreten Aufgaben. Zuerst müssen die Zustände und ihre Ladungsdarstellung gewählt werden. Dann braucht die Massenauslese eine native Operatoridentifikation, einen Auswahlmechanismus und eine Skala. Schließlich muss sich dieselbe Materiekonstruktion in einer erhaltenen Quelle niederschlagen, deren Wirkung sich geometrisch beschreiben und prüfen lässt. Die Rechnungen geben an, an welchen Übergängen diese Schritte voneinander abhängen.

\subsection{Drehsinn und Turm der Hilbert-Shifts}
\label{sec:handedness-tower}

Nach Fassung 0.2.0 sind zwei Fragen in einer Reihe von Erkundungen und zwei Tests weiterverfolgt worden; Runde v4.10.0-r59 des Companion
fasst sie mit ihrem Status zusammen \cite{towerlayer}. Bis auf die beiden Tests H125 und H126 sind alle Schritte Erkundungen, deren Ergebnis
vor der Rechnung absehbar war; sie sind nicht blind und prüfen nichts.

\paragraph{Woher der Drehsinn kommt.}
Abschnitt~\ref{sec:octant-sign} und Runde v4.10.0-r58 führen das Vorzeichen des Oktantschritts auf das Produkt \(\tau\varepsilon\) aus der
Zeitrichtung der Auslese und dem Umlaufsinn der Bahn zurück. Fünf Erkundungen haben gefragt, ob ein weiterer Baustein diesen Drehsinn
festlegt: der Nullpunkt der Fläche einer Punktierung, die Restdrehung einer Scheibe aus drei Bögen als Holonomie, das Licht als Vermittler
zwischen den Frames, eine Störung des Lichtgewichts und die Drehrate einer Scheibe als Ursprung von \(\Lambda\). Jede Rechnung setzt den
Drehsinn voraus, keine legt ihn fest. Zwei Nebenergebnisse sind exakt: Mit dem halbklassischen Spektrum \(j+\tfrac12\) statt
\(\sqrt{j(j+1)}\) wird die Zählung geschlossen lösbar und gibt genau den Hopf-Wert \(\beta_{\rm sym}=1/\sqrt3\), verfehlt aber die Daten;
und der größte Formzyklus der Drei-Bogen-Scheibe dreht sie um genau \(-\pi\), sodass die Restdrehung \(\gamma_M\cdot\tfrac12\) wird, wenn man
\(\gamma_M\) als Kopplung einsetzt, nicht herleitet. Die Verknüpfung von \(\Lambda\) mit einer Drehrate ist nur eine Umbenennung, wie die
Epoche in Runde v4.10.0-r56. Am Ende ist die offene Wahl bei \(\varepsilon\) angekommen, also bei der Wahl der imaginären Einheit; über
\(\mathbb C\) ist das eine Konvention.

\paragraph{Ein relatives Vorzeichen (H125).}
Messbar ist dagegen ein Drehsinn zweier Sektoren relativ zueinander. Test H125 wurde vor der Rechnung festgelegt. Seine erste Stufe, eine
Strukturrechnung mit vorab erwartetem Ausgang, ist bestätigt: Die Familien-\(su(2)\) der Sedenionen (Runde v4.10.0-r14.4) ist auf den
Leptonen relativ zum Lorentz-Spin das komplex konjugierte Dublett, auf den Quarks nicht; die Bargmann-Invarianten sind
\(\tfrac14-\tfrac i4\) und \(\tfrac14+\tfrac i4\), in allen \(24+24\) zulässigen Zuordnungen. Die Wahl \(i\mapsto-i\) konjugiert beide Sektoren
zugleich und lässt das relative Vorzeichen stehen. Die zweite Stufe ist blind gegenüber künftigen Daten: Unter der Setzung, dass die
CP-Phase in beiden Sektoren auf dieselbe Weise aus der Familienstruktur entsteht, folgt \(\operatorname{sign}J_{\rm PMNS}=-\operatorname{sign}J_{\rm CKM}\),
also \(\sin\delta_{\rm PMNS}<0\). Entschieden wird mit der ersten Auswertung, die das Vorzeichen auf mindestens \(3\sigma\) festlegt,
spätestens Ende 2035. Ein einzelnes Vorzeichen hat die Grundwahrscheinlichkeit \(\tfrac12\); eine Bestätigung stützte vor allem die Setzung.
Runde v4.10.0-r65 schränkt diesen Befund ein: Das relative Vorzeichen hängt an einer weiteren benannten Setzung
(Abschnitt~\ref{sec:family-copies-koide}).

\paragraph{Der Turm der Hilbert-Shifts.}
Die zweite Frage war, ob der Feldwert eine Erzeugungsregel von einem Hilbert-Shift zum nächsten hat. Liest man den Feldwert auf der
Stufe \(n\) als Element der Cayley--Dickson-Algebra der Dimension \(2^n\), mit der Grundbahn \(\psi(t)=\exp(u\omega t)\psi_0\), dann ist die erste
Stufe eine reelle Schwingung, die zweite die Ellipse mit der Fläche als \(U(1)\)-Ladung, die dritte die Hopf-Sphäre mit der ersten
Holonomie und \(SU(2)\), die vierte die Oktonionen mit \(G_2\) und \(su(3)\) als Stabilisator einer Einheit. Danach kommt keine stetige
Symmetrie hinzu; die Kreisbahn schließt nicht mehr für jede Richtung, und Nullteiler treten auf. Neu ist bei den Sedenionen eine diskrete
Symmetrie \(S_3\), die drei Oktonionen-Kopien vertauscht, und die Zahl verträglicher Achsen folgt der Bott-Periode. Die Algebra ist bekannt;
neu ist nur die Zuordnung zum Turm, und sie ist eine Lesart. Rückwärts gelesen sind die meisten Wahlen oberhalb der ersten Stufe erzwungen
oder Konvention: Im Turm-Modell ist unter den erreichbaren Farbzahlen nur \(N_c=3\) anomaliefrei, ein Vorzeichenwechsel in der
Grundgleichung führt auf gespaltene statt kompakte Symmetrien, und eine CP-Phase gibt es erst ab drei Generationen. Mit einem Faktor
\(\exp(u\pi/4)\) in Kopienform erfüllt ein Zustand die Koide-Bedingung genau; die Phase \(\delta\) bleibt frei, und drei vollständige
Generationen passen in diese Algebra nicht. Ein Test, die Neutrino-Mischung aus der nicht-assoziativen Kopplung zweier Oktonion-Richtungen
zu gewinnen (H126), ist nach dem vorab festgelegten Kriterium nicht prüfbar, weil die Konstruktion überall entartet. Ein nichtlinearer
Zusatz in der Grundgleichung bringt eine Skala, ändert die Struktur aber nicht.

Keine dieser Rechnungen leitet \(\beta\), eine Masse, einen Mischungswinkel oder \(\Lambda\) her. Gesetzt bleiben der Rahmen der
Chiralität, die Unterscheidung von M und T, die Zahl der Generationen und jede Skala. OP1, OP3, OP4 und OP5 bleiben offen.

\subsection{Familiensymmetrie, Kopien und der Koide-Punkt}
\label{sec:family-copies-koide}

Nach Fassung 0.3.0 hat der Companion die Runden v4.10.0-r60 bis v4.10.0-r67 erhalten \cite{layers}. Die Runden r66 und r67 betreffen die
Zahl \(\beta\) und stehen in Abschnitt~\ref{sec:one-number}. Die Runden r60 bis r65 sind Sammelrunden und eine Korrekturrunde: Sie rechnen
nichts neu und halten Tests und Erkundungen mit dem Status fest, den die Ergebnis-Notizen nennen \cite{roundsr60r65}. Dieser Abschnitt gibt
sie verkürzt wieder; maßgeblich sind die Runden. Die erste Tabelle nennt die Tests mit vorab festgelegtem Kriterium. Wo ein Teil nicht blind
war, steht das in der dritten Spalte; die Runden führen es je Teilaussage.

\begin{center}
\small
\begin{tabular}{@{}p{.21\linewidth}p{.09\linewidth}p{.15\linewidth}p{.45\linewidth}@{}}
\hline
\textbf{Test} & \textbf{Runde} & \textbf{blind?} & \textbf{Ergebnis}\\
\hline
\raggedright H127: Brechung der Familiensymmetrie durch die Ladung & r60.2 & \raggedright vorab festgelegt & \raggedright Stufe A bestätigt: Die Brechung \(SO(3)\to A_4\) ist ein reiner Spin-3-Tensor. Stufe B nach dem Wortlaut widerlegt; sie war für jede zulässige Kopplung leer (r60.3).\tabularnewline
\raggedright H138--H147: Oktonion-Kopien in den Sedenionen, Rest-Generation & r62 & \raggedright teils nicht blind & \raggedright Widerlegt: H138a, H138b, H142b, H144c, nach dem Wortlaut H139a und H145c; die übrigen Teile bestätigt. Es gibt genau drei \(G_2\)-stabile Oktonion-Kopien, bedingt an die Definition der Kopie. Sie sind drei Lesarten einer Generation, nicht drei Generationen. Die Kopienstruktur liefert keine Massen.\tabularnewline
\raggedright H148--H156: der Koide-Punkt & r63 & \raggedright teils nicht blind & \raggedright Widerlegt: H150b in einem Teil, H155c, nach dem Wortlaut H150c und H155d; teilweise H151b und H156c; die übrigen Teile bestätigt, H153b bis auf eine Teilvorhersage. \(K=\tfrac23\) folgt aus einem Postulat mit gesetzter Skala, nicht aus der Algebra.\tabularnewline
\raggedright H128, H130: Familiensymmetrie und Träger in den Sedenionen & r64.1, r64.2 & \raggedright teils nicht blind & \raggedright H128a widerlegt (Dimension 22 statt der vorab genannten 17, numerischer Rang); H128b für die Richtung \(e_7\) bestätigt, für \(e_8\) widerlegt; H130a und H130b bestätigt.\tabularnewline
\raggedright H131, H134: Mischung aus den Kopien & r64.3 & \raggedright teils nicht blind & \raggedright H131a bestätigt, H131b widerlegt; H134a und H134c bestätigt, H134b widerlegt. Verträglichkeit mit den Daten, keine Vorhersage.\tabularnewline
\raggedright H129, H132, H133, H135, H136: rekursive Feldgleichung auf dem Turm & r64.4 & \raggedright teils nicht blind & \raggedright H129 und H133 nicht bestätigt; H132, H135 und H136 widerlegt. Der Ansatz legt keine Skala und keinen Wert fest.\tabularnewline
\raggedright H105: Flächenschritte und Echos & r64.5 & \raggedright vorab festgelegt, gegen künftige Daten & \raggedright offen; Frist 31.~Dezember 2028.\tabularnewline
\raggedright H162: Ladungsachse auf dem Sedenionen-Träger & r65.1 bis r65.3 & \raggedright a, b, d nicht blind & \raggedright a bis d bestätigt, c nur numerisch; Teil e ist eine Erkundung.\tabularnewline
\hline
\end{tabular}
\end{center}

\paragraph{Familiensymmetrie und Kopien.}
Die hermiteschen \(3\times3\)-Matrizen über den Oktonionen bilden die Albert-Algebra. Dort ist der Kommutant von \(G_2\) in
\(F_4\) eine Familien-\(SO(3)\); die Ladung bricht sie zu genau \(A_4\) (Runde v4.10.0-r60, Erkundungen k275 und k277, nicht blind). Die Gruppe \(A_4\),
die frühere Runden als Symmetrie der drei Plätze gesetzt hatten, ist dort also ein Rest. Bei exakter \(A_4\) sind die Quarkmassen entartet,
und die Quarkmischung ist die Einheitsmatrix; eine Mischung verlangt, dass \(A_4\) selbst bricht. Die Ausrichtung des Vakuums hängt an freien
Koeffizienten eines Potentials, sodass ein blinder Test derzeit nicht möglich ist. Die drei Plätze als Generationen zu lesen ist eine Lesart.
Runde v4.10.0-r62 fragt, ob die Sedenionen die Zahl der Generationen erklären. Die Abschlussbedingung einer Kopie ist kubisch und hat genau
drei Lösungen. Die drei Kopien haben dieselbe Farbe und tragen je eine Generation, aber keine gemeinsame Ladung, und zusammen nur 15 von 24
Zuständen. Auf der nächsten Stufe des Turms tritt bei vier Kopien in 9 von 126 Auswahlen eine vollständige Rest-Generation auf; ob ihre
Zustände physikalische Zustände sind, ist nicht gezeigt. Die Zahl der Generationen bleibt eine Eingabe.

\paragraph{Der Koide-Punkt.}
Liest man die Massen an den Amplituden ab (eine Setzung aus Runde v4.10.0-r59), dann gilt \(K=1/(3\cos^2\theta)\) mit einem Mischwinkel
\(\theta\), also \(K=\tfrac23\) genau bei \(\theta=\pi/4\) (Runde v4.10.0-r63). Dieser Winkel ist der Fixpunkt eines Tauschs \(J\), den der
Companion als Tausch von Materie- und Tachyon-Frame liest; jedes tauschsymmetrische Potential ist dort stationär. Hergeleitet sind diese
beiden Aussagen. Dass der Fixpunkt angenommen wird, folgt nur aus dem Postulat, dass die beiden Frames gleichberechtigt sind, und mit
gesetzter Skala; der Autor hat das Postulat am 26.~November 2026 als Arbeitsannahme angenommen, hergeleitet ist es nicht. Der Tausch \(J\)
ist antilinear und weder die Tachyon-Abbildung \(z\mapsto iz\) noch die Zeitumkehr \(z\mapsto\bar z\) (Bemerkung v4.10.0-r64.7). Die Phase
\(\delta=2/9\) bleibt angepasst. Ohne gesetzte Skala verschiebt der Fluss den Wert abhängig vom Start. Der gemessene Wert
\(K=0{,}6666645\pm0{,}0000051\) liegt \(0{,}43\sigma\) unter \(\tfrac23\). Runde r63 führt daraus eine bedingte, offene Vorhersage
\(m_\tau=1776{,}969\)\,MeV, falls \(K=\tfrac23\) für die Polmassen exakt gilt; sie ist nicht die Vorhersage H102 aus
Abschnitt~\ref{sec:data}.

\paragraph{Ladungsachse und das relative Vorzeichen.}
Abschnitt~\ref{sec:handedness-tower} nennt das relative CP-Vorzeichen zwischen Leptonen und Quarks (H125). Es war an eine Ladungsachse
gebunden, die ein Platzhalter des Companion ist (CT-003); Test H130 stützt diese Achse nicht (Runde v4.10.0-r64). Test H162 zeigt, dass das
Spektrum der Ladung, die Konstruktion mit zwei Generationen, die Familien-\(su(2)\) und die erste Stufe von H125 für jede Achse gelten
(Runde v4.10.0-r65). Eine Erkundung derselben Runde (k361, nicht blind, nur numerisch) findet zwei Typen der Ladung, mit Links- oder mit
Rechtsmultiplikation, und schränkt den Befund aus Abschnitt~\ref{sec:handedness-tower} ein: Das relative Vorzeichen hängt nicht an der Wahl
\(i\mapsto-i\), nicht an der Achse und nicht am Ladungstyp, wohl aber an der Setzung, dass dieselbe Seite der Multiplikation die Familie in
beiden Sektoren trägt. Sie liegt nahe, folgt aber nicht aus Ladung und Farbe allein. Die Vorhersage \(\sin\delta_{\rm PMNS}<0\) bleibt offen
und hängt jetzt an zwei benannten Setzungen. Der Träger enthält zwei Generationen, nicht drei.

\paragraph{Korrekturen und Nachrechnung.}
Runde v4.10.0-r61 berichtigt zwei Formeln früherer Runden (Erkundung k320, nicht blind): die Normierung der Vakuumdichte,
\(\rho_\Lambda=\Lambda/\kappa\), und die Wurzel in \(\beta=\sqrt{(\gamma_L-1)/(\gamma_L+1)}\). Keine Zahl, kein Testergebnis und kein Status
ändert sich; dieser Artikel hat keine der beiden Formeln zitiert. Runde r64 präzisiert außerdem den Geltungsbereich einiger älterer Aussagen,
ebenfalls ohne Statusänderung. Nach der Annahme des Zählpostulats hat eine Erkundung (k358, kein Test; Runde v4.10.0-r68) die Zahlen, die von \(\beta\)
abhängen, mit dem exakten Wert nachgerechnet \cite{recalc}: Kein Urteil ändert sich; am schärfsten reagiert \(m_W\), dessen Abstand zu
PDG~2026 von \(-0{,}44\sigma\) beim zuvor geführten, aus gerundeten Eingaben gewonnenen Wert auf \(-0{,}13\sigma\) geht, und \(m_\tau\) aus
H102 bleibt bei \(+2{,}3\sigma\).

Keine dieser Runden leitet \(\beta\), eine Masse, einen Mischungswinkel oder \(\Lambda\) her. Gesetzt bleiben die Zahl der Generationen,
die Definition der Kopie, die Ladungsachse, die Seite der Multiplikation, die Auslese an den Amplituden und jede Skala; angepasst bleibt
\(\delta\). OP1 bis OP5, C5 und C6/T2 bleiben offen, und kein Guard ist aufgehoben.

\subsection{Was nach dieser Fassung offen bleibt}
\label{sec:open-points}

Der Companion führt seine offenen Probleme unter festen Namen. Ein Guard ist dort eine ausdrücklich gesperrte Aussage: Sie darf nicht
als bewiesen verwendet werden, solange ihr Nachweis fehlt. Die folgende Tabelle nennt die Probleme mit einer Zeile; keines ist
geschlossen, und kein Guard ist aufgehoben \cite{companion,layers}.

\begin{center}
\small
\begin{tabular}{@{}p{.12\linewidth}p{.55\linewidth}p{.27\linewidth}@{}}
\hline
\textbf{Name} & \textbf{Frage} & \textbf{im Artikel}\\
\hline
OP1 & Woher kommt die Zahl \(\beta=0{,}5494\), also warum zählen die Frames M und T die Horizontzustände nach verschiedenen Regeln? & Abschnitt~\ref{sec:one-number}\\
OP2 & Existiert der Dirac-artige Operator auf dem antisymmetrischen Sektor als selbstadjungierter Operator, ohne die zurückgezogene Riemann-Arbeit, und liegt sein Spektralversatz bei \(\beta\)? & nur als Gleichsetzung in Abschnitt~\ref{sec:one-number}\\
OP3 & Welche Übergänge der Hülle sind zulässig, woher kommen drei Generationen, und was legt Phase und Vorzeichen der Leptonphase fest? & Abschnitte~\ref{sec:charged-multiplicity}, \ref{sec:octant-sign}, \ref{sec:handedness-tower} und \ref{sec:family-copies-koide}\\
OP4 & Woher kommen die Mischungswinkel der Neutrinos und der Quarks? & Gleichsetzung in Abschnitt~\ref{sec:one-number}; relatives CP-Vorzeichen in den Abschnitten~\ref{sec:handedness-tower} und \ref{sec:family-copies-koide}\\
OP5 & Wie fügen sich Vakuumbilanz und kosmologische Konstante in den Rahmen, und welche Materiequelle treibt welches geometrische Feld? & Fallstudie in Abschnitt~\ref{sec:op5-coupling}, Anhang~\ref{sec:anhang-b-annahmen}\\
C5 & Kann die verbleibende Obstruktion in endlich vielen Schritten exakt null werden, statt nur unter eine Detektorschwelle zu fallen? & nicht behandelt\\
C6/T2 & Schließen alle Auslesungen zugleich, global und nicht nur lokal? & nicht behandelt; durch Guards gesperrt\\
\hline
\end{tabular}
\end{center}

Drei Fragen tragen die weitere Arbeit. Sie sind hier mit ihrem Stand zusammengefasst; die Herleitungen dazu stehen im Companion.

Wie sie zusammenhängen, ordnet Runde v4.10.0-r55 des Companion \cite{crossmap,layers}. Alle bedingten Zahlenwerte hängen an der einen Zahl aus
Abschnitt~\ref{sec:one-number}. Quer zu den einzelnen Problemen liegen zwei Lücken. Die erste ist die Auslese: Mehrere Stellen brauchen ein
Bezugssystem, eine Achse und eine Zeitorientierung, die der Companion nicht festlegt. Sie zerfällt in die Frage nach der Orientierung, die
beim Vorzeichen des Oktantschritts und bei der Zählung der Horizontzustände auftritt, und die Frage nach Bezugssystem und Achse, die die
Beobachterlesart des Materie-Frames und die gerichtete OP5-Auslese betrifft. Die zweite Lücke ist eine gemeinsame Dynamik mit eigener Skala,
die bisher fehlt. Daneben hat jedes Problem einen eigenen Rest, der von der Zahl nicht abhängt. Die erste der drei Fragen unten gehört zur
Orientierung, die zweite zu einem solchen Rest, die dritte zu Bezugssystem und Achse und zur fehlenden Dynamik.

Für einen dieser Reste liegt inzwischen ein Negativ-Schluss vor. Der eigene Rest von OP5 ist der Absolutwert der kosmologischen
Konstante. Runde v4.10.0-r56 zeigt für eine ausdrücklich genannte Klasse von Konstruktionen des Companion (Frame-Abbildungen und
Auslesungen, die Wirkungen W1/W2, den semiklassischen Weg, das OP5-Vergleichsmodell und die Kaskade), dass eine additive Konstante aus
jeder Bewegungsgleichung herausfällt und mit spurfreien Gleichungen eine Integrationskonstante ist, dass diese Klasse keine eigene Länge
enthält und dass die Auslesung \(\Omega_M=\beta^2\) für jeden Wert von \(\Lambda\) zu genau einer Epoche gilt; sie gibt nur
\(\Lambda=3(1-\beta^2)H^2\) an dieser Epoche \cite{lambdainput}. Der Wert von \(\Lambda\) ist dort also eine Eingabe, gleichwertig die Wahl
der Epoche \glqq heute\grqq{}. Das ist ein Ausschluss für den vorhandenen Bestand, nicht für jede Erweiterung; OP5 bleibt offen, und die
Vorhersage \(\Omega_M=\beta^2\) wartet auf die Messung (Abschnitt~\ref{sec:one-number}).
Der andere Weg wäre eine Dynamik mit \(\Omega_M=\beta^2\) als Attraktor. Die einzige Dynamik des Companion für das Verhältnis, die Kaskade,
in der Vakuum mit der Rate \(\beta^2\varepsilon H\) in Materie zerfällt, hat den Attraktor \(\varepsilon\beta^2/3\); gleich \(\beta^2\) ist er genau für
\(\varepsilon=3\). Gegen die BAO-Daten von DESI DR2 ist dieses Modell nach einem vorab festgelegten, nicht blinden Vergleich nicht
ausgeschlossen, verlangt aber heute \(\Omega_M\approx0{,}50\) und passt deutlich schlechter als \(\Lambda\)CDM (Runde v4.10.0-r57 \cite{layers}).
Der Wert \(\beta^2\) läge dann in der Zukunft; warum er heute gemessen wird, erklärt der Attraktor nicht.

\paragraph{Ein nativer physikalischer Operator.}
Die Massenauslese braucht einen Operator, der aus dem gemeinsamen Träger folgt und nicht nachträglich gewählt ist. Bisher ist er nicht
gefunden. Abschnitt~\ref{sec:octant-sign} zeigt, dass der vorhandene Kompensator auf den geschlossenen ungeraden Zuständen verschwindet
und dass kein Baustein des Companion das Vorzeichen des Oktantschritts festlegt; es ist eine Zeitorientierung der Auslese. Dieselbe
Orientierung tritt bei der Zählung der Horizontzustände wieder auf: Das Gewicht 2 der Materie-Regel verlangt eine Auslese, die die beiden
Orientierungen einer Punktierung unterscheidet, ohne eine auszuwählen (Runde v4.10.0-r53 \cite{layers}). Ein nativer Operator müsste beides zugleich
erklären, das Vorzeichen im Einzelfall und die Unterscheidung in der Zählung. Eine Auslese in drei zeitlich geordneten Fenstern, die
Bargmann-Invariante des Oktant-Tripels auf einer Punktierung vom Spin \(j\), leistet das nicht (Runde v4.10.0-r58 \cite{layers}, Anhang~\ref{app:new-temporal}):
Sie gibt beiden Orientierungen denselben Betrag und die Werte \(2^{-3j}e^{\pm ij\pi/2}\), unterscheidet sie also nur für \(j\), die keine geraden
ganzen Zahlen sind. Als Zählregel gelesen, mit Gewicht 1 auf den geraden \(j\), ergibt sie \(\gamma_L'=1{,}838\) statt \(1{,}865\) und verfehlt
\(m_W\) um rund \(20\sigma\). Vorzeichen und Zählung kommen damit nicht aus derselben Auslese dieser Form. Abschnitt
\ref{sec:one-number} zeigt, dass alle bedingten Zahlenwerte an einer Zahl hängen, die über gesetzte Gleichsetzungen eingeht.

\paragraph{Eine hergeleitete Zuordnung der Eichladungen.}
Die Ladungszuordnung im Dixon-Vergleich ist konstruiert, nicht abgeleitet. Abschnitt~\ref{sec:charged-multiplicity} grenzt ein, wo ein
Generationsindex sitzen kann: nicht in der Farbe, sondern in einer farbneutralen Multiplizität, deren Zahl \(n\) eine Eingabe bleibt.

\paragraph{Ein unabhängig begründetes geometrisches Quellgesetz.}
Für den gewöhnlichen Teil ist es vorhanden: Der präparierte Beitrag in Abschnitt~\ref{sec:op5-coupling} folgt allein aus der
linearisierten Einstein-Gleichung und ist außerhalb seines Trägers kompakt (Anhang~\ref{sec:anhang-b-kompakt}). Der gerichtete Kontrast
braucht zusätzlich die Gleichsetzung der Higgsquelle mit dem Antrieb des Vergleichsfelds; am Higgs-Sektor der Wirkung W2 folgt sie nur für
den skalaren Kanal, und dann fehlt der gerichtete Vermittler, an dem der Kontrast hängt (Runde v4.10.0-r33 \cite{layers}). Die Suche nach
diesem Vermittler in den Runden v4.10.0-r34 bis v4.10.0-r54 trägt nur teilweise: Gesetzt bleiben ein massiver Spin-2-Vermittler, seine
Kopplung an den Higgsradius und eine aus Gründen der Kausalität verlangsamte Ausbreitung, und Bezugssystem und Achse sind bei OP5 eine
Setzung des Detektors, keine Wahl des Feldes \cite{layers}. Mit dem lichtschnellen Higgsradius von W2 ist die gerichtete Kopplung für keine
Wahl der Vermittler kausal verträglich, und im Vakuum begrenzen Laborgrenzen der Lorentz-Verletzung sie auf \(g_m^2\lesssim10^{-26}\); ein
beobachtbarer Kontrast verlangt ein Material. Offen ist damit, ob es einen solchen Vermittler als eigenes Feld gibt; die einzelnen Schritte
stehen in Anhang~\ref{sec:anhang-b-vermittler}.

Keiner dieser Punkte ist durch die vorliegende Fassung geschlossen. Sie legt offen, an welcher Stelle eine Herleitung ansetzen müsste.



\clearpage
\begingroup
\small
\raggedright
\begin{thebibliography}{99}
\setlength{\itemsep}{0.25\baselineskip}
\bibitem{companion}
M.~Brendecke, \emph{Quantum Sphaera Companion: A Structure-First Mathematical Unfolding}, v5.3.0, Vollausgabe, 2026.
\href{https://doi.org/10.5281/zenodo.24434403}{doi:10.5281/zenodo.24434403}.
\bibitem{physics}
M.~Brendecke, Quantum Sphaera Companion v5.3.0, Physik-Ausgabe, 2026.
\href{https://doi.org/10.5281/zenodo.24434403}{Gemeinsame Zenodo-Veröffentlichung}.
\bibitem{repo}
M.~Brendecke, \emph{QuantumSphere-Unfold}, Quellrepository.
\href{https://github.com/Moriarty1982/QuantumSphere-Unfold}{GitHub}.
\href{https://github.com/Moriarty1982/QuantumSphere-Unfold/tree/v5.3.0}{Release-Tag v5.3.0}.
\bibitem{grundidee}
M.~Brendecke, \emph{Mit dem heutigen Wissen rückwärts bis zur Grundidee}, Forschungsnotiz, 6. November 2026.
\href{https://github.com/Moriarty1982/QuantumSphere-Unfold/blob/f2a4e687273f2ec7f4124d490e696e07c00f387a/notes/research/qsc_back_to_basic_idea_2026_11_06.md}{Feldlesart und Frame-Zuordnung}.
\bibitem{feldkern}
M.~Brendecke, \href{https://github.com/Moriarty1982/QuantumSphere-Unfold/blob/f2a4e687273f2ec7f4124d490e696e07c00f387a/paper/short/p2_kernel.tex}{Kern und Lorentzstruktur}, Kurzentwurf zum Companion.
\bibitem{feldsektoren}
M.~Brendecke, \href{https://github.com/Moriarty1982/QuantumSphere-Unfold/blob/f2a4e687273f2ec7f4124d490e696e07c00f387a/paper/short/p3_frames.tex}{Dualitätssektoren und Kompensator}, Kurzentwurf zum Companion.
\bibitem{robertson}
H.~P. Robertson, \emph{The Uncertainty Principle}, Phys. Rev. 34, 163--164 (1929).
\href{https://doi.org/10.1103/PhysRev.34.163}{doi:10.1103/PhysRev.34.163}.

\bibitem{qam}
\emph{QAM-Dokumentation im Companion-Quellrepository}.
\href{https://github.com/Moriarty1982/QuantumSphere-Unfold/blob/f2a4e687273f2ec7f4124d490e696e07c00f387a/qam/README.md}{QAM und seine Prüfwerkzeuge}.
\bibitem{op5-density}
M.~Brendecke, \href{https://github.com/Moriarty1982/QuantumSphere-Unfold/blob/f2a4e687273f2ec7f4124d490e696e07c00f387a/notes/research/op5_bifocal_response_2026_11_18/OP5_J_W2_Dichteantrieb.md}{W2-Dichteantrieb und sein gravitativer Begleitterm}, OP5-Forschungsnotiz, 20. November 2026.
\bibitem{op5-channel}
M.~Brendecke, \href{https://github.com/Moriarty1982/QuantumSphere-Unfold/blob/f2a4e687273f2ec7f4124d490e696e07c00f387a/notes/research/op5_bifocal_response_2026_11_18/OP5_J_Luminaler_Einstein_Kanal.md}{Luminaler Einstein-Kanal und gerichteter massiver Kontrast}, OP5-Forschungsnotiz, 20. November 2026.
\bibitem{op5-motion}
M.~Brendecke, \href{https://github.com/Moriarty1982/QuantumSphere-Unfold/blob/f2a4e687273f2ec7f4124d490e696e07c00f387a/notes/research/op5_bifocal_response_2026_11_18/OP5_J_Dynamische_Gegenbewegung.md}{Dynamische Gegenbewegung mit kompakter linearer Gravitation}, OP5-Forschungsnotiz, 20. November 2026.
\bibitem{op5-axis}
M.~Brendecke, \href{https://github.com/Moriarty1982/QuantumSphere-Unfold/blob/f2a4e687273f2ec7f4124d490e696e07c00f387a/notes/research/op5_bifocal_response_2026_11_18/OP5_J_Gleicher_Messweg_Nullprofil.md}{Nullprofil und Krümmungsschaltung im selben Messweg}, OP5-Forschungsnotiz, 20. November 2026.
\bibitem{op5-elastic-limit}
M.~Brendecke, \href{https://github.com/Moriarty1982/QuantumSphere-Unfold/blob/f2a4e687273f2ec7f4124d490e696e07c00f387a/notes/research/op5_bifocal_response_2026_11_18/OP5_J_Hyperelastische_Halterung.md}{Hyperelastische Halterung in zweiter Ordnung}, OP5-Forschungsnotiz, 20. November 2026.
\bibitem{op5-reservoir}
M.~Brendecke, \href{https://github.com/Moriarty1982/QuantumSphere-Unfold/blob/f2a4e687273f2ec7f4124d490e696e07c00f387a/notes/research/op5_bifocal_response_2026_11_18/OP5_J_Reservoir_Quadrupolabgleich.md}{Reservoirabgleich der direkten Dichtebeiträge}, OP5-Forschungsnotiz, 20. November 2026.
\bibitem{dixon-commutant}
M.~Brendecke, \href{https://github.com/Moriarty1982/QuantumSphere-Unfold/blob/f2a4e687273f2ec7f4124d490e696e07c00f387a/sources/research/parts/v4100-post-release-layers.tex#L2582}{Kommutant einer Generation}, Companion, Runde v4.10.0-r19.
\bibitem{dixon-copies}
M.~Brendecke, \href{https://github.com/Moriarty1982/QuantumSphere-Unfold/blob/f2a4e687273f2ec7f4124d490e696e07c00f387a/notes/research/op3_dixon_copies_2026_11_16/README.md}{Zwei und vier Dixon-Kopien}, OP3-Forschungsnotiz, 16. November 2026.
\bibitem{op5-budget}
M.~Brendecke, \href{https://github.com/Moriarty1982/QuantumSphere-Unfold/blob/f2a4e687273f2ec7f4124d490e696e07c00f387a/notes/research/op5_bifocal_response_2026_11_18/OP5_J_AB_Fehlerbudget.md}{A/B-Fehlerbudget}, OP5-Forschungsnotiz.
\bibitem{op5-nullpfad}
M.~Brendecke, \href{https://github.com/Moriarty1982/QuantumSphere-Unfold/blob/f2a4e687273f2ec7f4124d490e696e07c00f387a/notes/research/op5_bifocal_response_2026_11_18/OP5_J_Nullpfad_Abgleich.md}{Nullpfad-Abgleich}, OP5-Forschungsnotiz.
\bibitem{op5-feder}
M.~Brendecke, \href{https://github.com/Moriarty1982/QuantumSphere-Unfold/blob/f2a4e687273f2ec7f4124d490e696e07c00f387a/notes/research/op5_bifocal_response_2026_11_18/OP5_J_Feder_Reservoirbilanz.md}{Federarbeit und Reservoirbilanz}, OP5-Forschungsnotiz.
\bibitem{op5}
M.~Brendecke, \emph{OP5-Forschungsübergabe nach PR 211}, 20. November 2026.
\href{https://github.com/Moriarty1982/QuantumSphere-Unfold/blob/f2a4e687273f2ec7f4124d490e696e07c00f387a/llm/Uebergabe.md}{GitHub-Forschungsstand}.
\bibitem{h115}
M.~Brendecke, \href{https://github.com/Moriarty1982/QuantumSphere-Unfold/blob/1368ae1febea0ce1a497878ddeb84060fc96af55/llm/Ergebnis_H115_Kompensator_Oktant.md}{H115: Kompensator und Nullrichtung des flachen Oktanten}, Testergebnis, 22. November 2026.
\bibitem{h116}
M.~Brendecke, \href{https://github.com/Moriarty1982/QuantumSphere-Unfold/blob/1368ae1febea0ce1a497878ddeb84060fc96af55/llm/Ergebnis_H116_Spinorband_Phasenquelle.md}{H116: Phasenauslese und Rückkehrquelle des Spinorbands}, Testergebnis, 22. November 2026.
\bibitem{h117}
M.~Brendecke, \href{https://github.com/Moriarty1982/QuantumSphere-Unfold/blob/1368ae1febea0ce1a497878ddeb84060fc96af55/llm/Ergebnis_H117_Zwei_Umlaeufe.md}{H117: Zwei Umläufe und die Wahl der Oktant-Orientierung}, Testergebnis, 22. November 2026.
\bibitem{rootmap}
M.~Brendecke, \href{https://github.com/Moriarty1982/QuantumSphere-Unfold/blob/1368ae1febea0ce1a497878ddeb84060fc96af55/notes/research/op_root_map_2026_11_18/README.md}{Wurzelkarte der offenen Probleme}, Forschungsnotiz, 18. November 2026.
\bibitem{chance}
M.~Brendecke, \href{https://github.com/Moriarty1982/QuantumSphere-Unfold/blob/1368ae1febea0ce1a497878ddeb84060fc96af55/notes/research/chance_estimate_2026_11_18/README.md}{Zufallsabschätzung der Treffer von \(\gamma_L\)}, Forschungsnotiz, 18. November 2026.
\bibitem{rovelli}
C.~Rovelli und S.~Speziale, \emph{Reconcile Planck-scale discreteness and the Lorentz-Fitzgerald contraction}, Phys. Rev. D 67, 064019 (2003).
\href{https://doi.org/10.1103/PhysRevD.67.064019}{doi:10.1103/PhysRevD.67.064019}.
\end{thebibliography}
\endgroup
\clearpage
\appendix
\newgeometry{hmargin=16mm,vmargin=19mm}
\section{Wolfram-Eingaben und Originalausgaben zur Fassung 0.1.0}
\label{app:wolfram-output}

Die Fenster zeigen links das Eingabeskript und rechts die zugehörige Originalausgabe. Die Zeilenangaben beziehen sich auf die jeweiligen Dateien. Die Blöcke eines Laufs setzen den vorangehenden Wolfram-Zustand voraus. Lange Zeilen werden für den Satz umgebrochen.

\medskip
\noindent
\begin{tabular}{@{}lp{.59\linewidth}r@{}}
\hline
\textbf{Lauf} & \textbf{Rechnungen} & \textbf{Seite} \\
\hline
HC und Feldalgebra & Projektionen, Biquaternionen und geometrische Identitäten & \pageref{app:wl-core} \\
Dixon-Träger & Kommutanten und Hyperladung & \pageref{app:wl-dixon} \\
Koide & Wurzeltripel, Gram-Spektrum und Massenauslese & \pageref{app:wl-koide} \\
OP5 & Quellen, Achsenprofil und Vergleichsrechnungen & \pageref{app:wl-op5} \\
\hline
\end{tabular}

\newcommand{\wolframdatei}[3]{%
  \ifnum\pdfstrcmp{#1}{dixon/dixon.wl}=0
    \VerbatimInput[firstline=#2,lastline=#3,fontsize=\footnotesize,baselinestretch=.88,%
      breaklines,breakafter={,[-/}]{audit/wolfram-v0.1.0/#1}%
  \else
    \VerbatimInput[firstline=#2,lastline=#3,fontsize=\footnotesize,baselinestretch=.88,%
      breaklines,breakanywhere]{audit/wolfram-v0.1.0/#1}%
  \fi}
\newcommand{\wolframfenster}[7]{%
  \subsubsection*{#1}%
  \noindent
  \begin{tabular}{|p{\dimexpr.5\linewidth-2\tabcolsep-1.5\arrayrulewidth\relax}%
                  |p{\dimexpr.5\linewidth-2\tabcolsep-1.5\arrayrulewidth\relax}|}
  \hline
  \rule{0pt}{2.5ex}\textbf{Wolfram-Input} \hfill \footnotesize Zeilen #3--#4 &
  \rule{0pt}{2.5ex}\textbf{Originalausgabe} \hfill \footnotesize Zeilen #6--#7 \\[.35ex]
  \hline
  \wolframdatei{#2}{#3}{#4}\vspace{.7\baselineskip} &
  \wolframdatei{#5}{#6}{#7}\vspace{.7\baselineskip} \\
  \hline
  \end{tabular}\par}

\medskip
\subsection{Heisenberg-Kompensator und Feldalgebra}\label{app:wl-core}
\wolframfenster{Projektion und ebene Welle}{core/article_core.wl}{1}{30}{core/article_core.txt}{1}{4}

\clearpage
\wolframfenster{Feldinvarianten, Biquaternionen und Nullspinoren}{core/article_core.wl}{31}{59}{core/article_core.txt}{5}{11}

\clearpage
\wolframfenster{Abschlussbedingung und Rang-eins-Vergleich}{core/article_core.wl}{60}{76}{core/article_core.txt}{12}{14}

\clearpage
\wolframfenster{Auslese, Eichung, Quellenreduktion und Koframe}{core/article_core.wl}{77}{102}{core/article_core.txt}{15}{23}

\clearpage
\subsection{Dixon-Träger und Kommutanten}\label{app:wl-dixon}
\wolframfenster{Multiplikation, Oktonionenderivationen und Farbkommutant}{dixon/dixon.wl}{1}{35}{dixon/dixon.txt}{1}{3}

\clearpage
\wolframfenster{Schwache Generatoren und Hyperladung}{dixon/dixon.wl}{36}{50}{dixon/dixon.txt}{4}{6}

\clearpage
\wolframfenster{Linearer Kommutant, Spektrum und antilineare Bedingung}{dixon/dixon.wl}{51}{77}{dixon/dixon.txt}{7}{16}

\clearpage
\subsection{Koide und Massenauslese}\label{app:wl-koide}
\wolframfenster{Zerlegung des Wurzeltripels und Phasenbereich}{koide/koide.wl}{1}{21}{koide/koide.txt}{1}{7}

\clearpage
\wolframfenster{HC-Projektion und Koisometrie}{koide/koide.wl}{22}{31}{koide/koide.txt}{8}{11}

\clearpage
\wolframfenster{Gram-Spektrum, Bargmann-Phase und Oktant}{koide/koide.wl}{32}{51}{koide/koide.txt}{12}{21}

\clearpage
\wolframfenster{Mischung, Interferenz und Massenskala}{koide/koide.wl}{52}{68}{koide/koide.txt}{22}{34}

\clearpage
\subsection{Geometrischer Vergleich in OP5}\label{app:wl-op5}
\wolframfenster{Quellterme und gekoppelte elastische Gleichungen}{op5/op5.wl}{1}{29}{op5/op5-output.txt}{1}{8}

\clearpage
\wolframfenster{Lineare Bilanz, Quellkontraktionen und Feldantwort}{op5/op5.wl}{30}{50}{op5/op5-output.txt}{9}{18}

\clearpage
\wolframfenster{Transversales Profil, Achsenkontrast und Kernintegral}{op5/op5.wl}{51}{65}{op5/op5-output.txt}{19}{24}

\clearpage
\wolframfenster{Quadratische Terme und Kontrollen des Quotienten}{op5/op5.wl}{66}{87}{op5/op5-output.txt}{25}{36}


\end{document}
